\documentclass[11pt,a4paper]{article}
\usepackage[a4paper,margin=20mm]{geometry}
\usepackage{fontspec}
\setmainfont{Noto Serif}
\setsansfont{Noto Sans}
\usepackage{amsmath,amssymb}
\usepackage{booktabs,tabularx,array}
\usepackage[table]{xcolor}
\usepackage{hyperref}
\definecolor{primary}{RGB}{27,79,114}
\definecolor{accent}{RGB}{40,116,166}
\definecolor{rowtint}{RGB}{235,245,251}
\hypersetup{colorlinks=true,linkcolor=primary,urlcolor=primary}
\setlength{\parskip}{5pt plus 2pt minus 1pt}
\setlength{\parindent}{0pt}
\begin{document}
\begin{center}
  {\LARGE\bfseries\color{primary} PHIẾU HỌC TẬP TRẮC NGHIỆM TOÁN 12}\\[4pt]
  {\normalsize\itshape\color{accent} Bộ sách: Kết nối tri thức với cuộc sống (Tập 1 \& Tập 2 – 6 Chương)}\\[4pt]
  {\small Cấu trúc mỗi chương: 20 câu Dễ (Câu 1--20) $\bullet$ 20 câu Trung bình (Câu 21--40) $\bullet$ 10 câu Khó (Câu 41--50) -- Tổng 50 câu/chương}
\end{center}
\vspace{6pt}\hrule\vspace{10pt}
\section*{\color{primary} CHƯƠNG I: ỨNG DỤNG ĐẠO HÀM ĐỂ KHẢO SÁT VÀ VẼ ĐỒ THỊ HÀM SỐ}
\textit{Trọng tâm kiến thức SGK: Tính đơn điệu và cực trị của hàm số, giá trị lớn nhất và giá trị nhỏ nhất của hàm số, đường tiệm cận của đồ thị hàm số (đứng, ngang, xiên) và khảo sát sự biến thiên của hàm số.}\par\vspace{4pt}
\subsection*{\color{accent} Phần I. Mức độ Dễ (Nhận biết -- Câu 1 đến Câu 20)}
\textbf{\color{primary} Câu 1.} Hàm số bậc ba y = x$^3$ - 3x$^2$ + 2 nghịch biến trên khoảng nào dưới đây?\par
\hspace*{1.5em}\textbf{A.} (0; 2)\quad\quad \textbf{B.} (-2; 0)\quad\quad \textbf{C.} (2; +$\infty$)\quad\quad \textbf{D.} (-$\infty$; 0)\par\vspace{2pt}
\textbf{\color{primary} Câu 2.} Thực hiện yêu cầu nhận biết sau: Phương trình đường tiệm cận đứng và tiệm cận ngang của đồ thị hàm số y = (5x - 1)/(x - 3) lần lượt là.\par
\hspace*{1.5em}\textbf{A.} x = 5 và y = 3\quad\quad \textbf{B.} x = -3 và y = 5\quad\quad \textbf{C.} x = 3 và y = 5\quad\quad \textbf{D.} x = 3 và y = -1\par\vspace{2pt}
\textbf{\color{primary} Câu 3.} Trong bài học thuộc Chương I, xét nội dung: Điểm cực đại của hàm số y = -x$^3$ + 3x + 6 là. Đáp án đúng là:\par
\hspace*{1.5em}\textbf{A.} x = 0\quad\quad \textbf{B.} x = -1\quad\quad \textbf{C.} x = 1\quad\quad \textbf{D.} y = 8\par\vspace{2pt}
\textbf{\color{primary} Câu 4.} Dựa vào kiến thức chuẩn SGK của Chương I (Ứng dụng đạo hàm để khảo sát và vẽ đồ thị hàm số), hàm số bậc ba y = x$^3$ - 3x$^2$ + 5 nghịch biến trên khoảng nào dưới đây:\par
\hspace*{1.5em}\textbf{A.} (-2; 0)\quad\quad \textbf{B.} (2; +$\infty$)\quad\quad \textbf{C.} (0; 2)\quad\quad \textbf{D.} (-$\infty$; 0)\par\vspace{2pt}
\textbf{\color{primary} Câu 5.} Chọn kết quả đúng nhất khi xét: Phương trình đường tiệm cận đứng và tiệm cận ngang của đồ thị hàm số y = (11x - 1)/(x - 6) lần lượt là:\par
\hspace*{1.5em}\textbf{A.} x = 6 và y = -1\quad\quad \textbf{B.} x = 11 và y = 6\quad\quad \textbf{C.} x = -6 và y = 11\quad\quad \textbf{D.} x = 6 và y = 11\par\vspace{2pt}
\textbf{\color{primary} Câu 6.} Khi làm bài tập cơ bản về ứng dụng đạo hàm để khảo sát và vẽ đồ thị hàm số, một học sinh gặp câu hỏi: Điểm cực đại của hàm số y = -x$^3$ + 3x + 9 là. Kết quả nào dưới đây đúng?\par
\hspace*{1.5em}\textbf{A.} x = -1\quad\quad \textbf{B.} x = 1\quad\quad \textbf{C.} y = 11\quad\quad \textbf{D.} x = 0\par\vspace{2pt}
\textbf{\color{primary} Câu 7.} Xác định giá trị hoặc khẳng định đúng cho bài toán sau: Hàm số bậc ba y = x$^3$ - 3x$^2$ + 8 nghịch biến trên khoảng nào dưới đây:\par
\hspace*{1.5em}\textbf{A.} (-2; 0)\quad\quad \textbf{B.} (-$\infty$; 0)\quad\quad \textbf{C.} (0; 2)\quad\quad \textbf{D.} (2; +$\infty$)\par\vspace{2pt}
\textbf{\color{primary} Câu 8.} Áp dụng trực tiếp định nghĩa và quy tắc trong Chương I: Phương trình đường tiệm cận đứng và tiệm cận ngang của đồ thị hàm số y = (17x - 1)/(x - 9) lần lượt là:\par
\hspace*{1.5em}\textbf{A.} x = 9 và y = -1\quad\quad \textbf{B.} x = 9 và y = 17\quad\quad \textbf{C.} x = 17 và y = 9\quad\quad \textbf{D.} x = -9 và y = 17\par\vspace{2pt}
\textbf{\color{primary} Câu 9.} Kiểm tra nhanh kiến thức nền tảng: Điểm cực đại của hàm số y = -x$^3$ + 3x + 12 là:\par
\hspace*{1.5em}\textbf{A.} x = -1\quad\quad \textbf{B.} x = 1\quad\quad \textbf{C.} y = 14\quad\quad \textbf{D.} x = 0\par\vspace{2pt}
\textbf{\color{primary} Câu 10.} Cho dữ kiện nhận biết sau đây: Hàm số bậc ba y = x$^3$ - 3x$^2$ + 11 nghịch biến trên khoảng nào dưới đây. Phương án chính xác là:\par
\hspace*{1.5em}\textbf{A.} (-2; 0)\quad\quad \textbf{B.} (2; +$\infty$)\quad\quad \textbf{C.} (-$\infty$; 0)\quad\quad \textbf{D.} (0; 2)\par\vspace{2pt}
\textbf{\color{primary} Câu 11.} Hãy chỉ ra kết quả đúng trong bốn lựa chọn A, B, C, D đối với yêu cầu: Phương trình đường tiệm cận đứng và tiệm cận ngang của đồ thị hàm số y = (23x - 1)/(x - 12) lần lượt là:\par
\hspace*{1.5em}\textbf{A.} x = 12 và y = 23\quad\quad \textbf{B.} x = 23 và y = 12\quad\quad \textbf{C.} x = -12 và y = 23\quad\quad \textbf{D.} x = 12 và y = -1\par\vspace{2pt}
\textbf{\color{primary} Câu 12.} Câu hỏi ôn tập lý thuyết và kỹ năng cơ bản (Chương I): Điểm cực đại của hàm số y = -x$^3$ + 3x + 15 là:\par
\hspace*{1.5em}\textbf{A.} x = -1\quad\quad \textbf{B.} x = 0\quad\quad \textbf{C.} x = 1\quad\quad \textbf{D.} y = 17\par\vspace{2pt}
\textbf{\color{primary} Câu 13.} Theo nội dung sách giáo khoa Kết nối tri thức với cuộc sống (Chương I), hàm số bậc ba y = x$^3$ - 3x$^2$ + 14 nghịch biến trên khoảng nào dưới đây:\par
\hspace*{1.5em}\textbf{A.} (0; 2)\quad\quad \textbf{B.} (-2; 0)\quad\quad \textbf{C.} (-$\infty$; 0)\quad\quad \textbf{D.} (2; +$\infty$)\par\vspace{2pt}
\textbf{\color{primary} Câu 14.} Tìm đáp số đúng của bài toán cơ bản sau: Phương trình đường tiệm cận đứng và tiệm cận ngang của đồ thị hàm số y = (29x - 1)/(x - 15) lần lượt là:\par
\hspace*{1.5em}\textbf{A.} x = 15 và y = -1\quad\quad \textbf{B.} x = 15 và y = 29\quad\quad \textbf{C.} x = -15 và y = 29\quad\quad \textbf{D.} x = 29 và y = 15\par\vspace{2pt}
\textbf{\color{primary} Câu 15.} Xét đối tượng toán học trong Chương I: Điểm cực đại của hàm số y = -x$^3$ + 3x + 18 là. Giá trị hoặc tính chất tương ứng là:\par
\hspace*{1.5em}\textbf{A.} y = 20\quad\quad \textbf{B.} x = -1\quad\quad \textbf{C.} x = 0\quad\quad \textbf{D.} x = 1\par\vspace{2pt}
\textbf{\color{primary} Câu 16.} Hoàn thành phép xác định cơ bản sau đây: Hàm số bậc ba y = x$^3$ - 3x$^2$ + 17 nghịch biến trên khoảng nào dưới đây:\par
\hspace*{1.5em}\textbf{A.} (0; 2)\quad\quad \textbf{B.} (-2; 0)\quad\quad \textbf{C.} (2; +$\infty$)\quad\quad \textbf{D.} (-$\infty$; 0)\par\vspace{2pt}
\textbf{\color{primary} Câu 17.} Trong các khẳng định hoặc giá trị dưới đây, đâu là đáp án đúng cho câu hỏi: Phương trình đường tiệm cận đứng và tiệm cận ngang của đồ thị hàm số y = (35x - 1)/(x - 18) lần lượt là?\par
\hspace*{1.5em}\textbf{A.} x = -18 và y = 35\quad\quad \textbf{B.} x = 35 và y = 18\quad\quad \textbf{C.} x = 18 và y = 35\quad\quad \textbf{D.} x = 18 và y = -1\par\vspace{2pt}
\textbf{\color{primary} Câu 18.} Bài tập khởi động Chương I: Điểm cực đại của hàm số y = -x$^3$ + 3x + 21 là:\par
\hspace*{1.5em}\textbf{A.} x = -1\quad\quad \textbf{B.} x = 0\quad\quad \textbf{C.} x = 1\quad\quad \textbf{D.} y = 23\par\vspace{2pt}
\textbf{\color{primary} Câu 19.} Quan sát dữ kiện đề bài: Hàm số bậc ba y = x$^3$ - 3x$^2$ + 20 nghịch biến trên khoảng nào dưới đây. Hãy chọn phương án đúng:\par
\hspace*{1.5em}\textbf{A.} (-$\infty$; 0)\quad\quad \textbf{B.} (-2; 0)\quad\quad \textbf{C.} (0; 2)\quad\quad \textbf{D.} (2; +$\infty$)\par\vspace{2pt}
\textbf{\color{primary} Câu 20.} Sử dụng công thức hoặc tính chất đầu bài của Chương I để giải quyết yêu cầu: Phương trình đường tiệm cận đứng và tiệm cận ngang của đồ thị hàm số y = (41x - 1)/(x - 21) lần lượt là:\par
\hspace*{1.5em}\textbf{A.} x = 21 và y = -1\quad\quad \textbf{B.} x = 41 và y = 21\quad\quad \textbf{C.} x = -21 và y = 41\quad\quad \textbf{D.} x = 21 và y = 41\par\vspace{2pt}
\subsection*{\color{accent} Phần II. Mức độ Trung bình (Thông hiểu -- Câu 21 đến Câu 40)}
\textbf{\color{primary} Câu 21.} Giá trị lớn nhất của hàm số f(x) = x$^3$ - 3x + 5 trên đoạn [0; 2] bằng:\par
\hspace*{1.5em}\textbf{A.} 3\quad\quad \textbf{B.} 9\quad\quad \textbf{C.} 5\quad\quad \textbf{D.} 7\par\vspace{2pt}
\textbf{\color{primary} Câu 22.} Vận dụng quy tắc trong Chương I (Ứng dụng đạo hàm để khảo sát và vẽ đồ thị hàm số) để giải bài toán: Đường tiệm cận xiên của đồ thị hàm số phân thức y = x + 3 + 3/(x - 1) có phương trình là:\par
\hspace*{1.5em}\textbf{A.} y = x + 3\quad\quad \textbf{B.} y = x - 1\quad\quad \textbf{C.} x = 1\quad\quad \textbf{D.} y = x + 6\par\vspace{2pt}
\textbf{\color{primary} Câu 23.} Biến đổi và tính toán theo kiến thức Chương I: Tổng giá trị cực đại và giá trị cực tiểu của hàm số y = x$^3$ - 3x + 6 bằng. Kết quả thu được là:\par
\hspace*{1.5em}\textbf{A.} 4\quad\quad \textbf{B.} 0\quad\quad \textbf{C.} 12\quad\quad \textbf{D.} 6\par\vspace{2pt}
\textbf{\color{primary} Câu 24.} Xét bài toán mức độ thông hiểu sau: Giá trị lớn nhất của hàm số f(x) = x$^3$ - 3x + 8 trên đoạn [0; 2] bằng. Đáp án đúng là:\par
\hspace*{1.5em}\textbf{A.} 10\quad\quad \textbf{B.} 8\quad\quad \textbf{C.} 6\quad\quad \textbf{D.} 12\par\vspace{2pt}
\textbf{\color{primary} Câu 25.} Để hoàn thành phiếu học tập Chương I, hãy xác định: Đường tiệm cận xiên của đồ thị hàm số phân thức y = x + 6 + 3/(x - 1) có phương trình là:\par
\hspace*{1.5em}\textbf{A.} y = x - 1\quad\quad \textbf{B.} y = x + 9\quad\quad \textbf{C.} x = 1\quad\quad \textbf{D.} y = x + 6\par\vspace{2pt}
\textbf{\color{primary} Câu 26.} Phân tích dữ kiện bài toán: Tổng giá trị cực đại và giá trị cực tiểu của hàm số y = x$^3$ - 3x + 9 bằng. Giá trị cần tìm bằng:\par
\hspace*{1.5em}\textbf{A.} 18\quad\quad \textbf{B.} 9\quad\quad \textbf{C.} 0\quad\quad \textbf{D.} 4\par\vspace{2pt}
\textbf{\color{primary} Câu 27.} Áp dụng tính chất đã học ở Chương I vào bài tập sau: Giá trị lớn nhất của hàm số f(x) = x$^3$ - 3x + 11 trên đoạn [0; 2] bằng:\par
\hspace*{1.5em}\textbf{A.} 13\quad\quad \textbf{B.} 15\quad\quad \textbf{C.} 9\quad\quad \textbf{D.} 11\par\vspace{2pt}
\textbf{\color{primary} Câu 28.} Thực hiện các bước suy luận hợp lý để giải quyết yêu cầu: Đường tiệm cận xiên của đồ thị hàm số phân thức y = x + 9 + 3/(x - 1) có phương trình là:\par
\hspace*{1.5em}\textbf{A.} y = x + 12\quad\quad \textbf{B.} y = x - 1\quad\quad \textbf{C.} x = 1\quad\quad \textbf{D.} y = x + 9\par\vspace{2pt}
\textbf{\color{primary} Câu 29.} Cho bài toán tự luyện (Ứng dụng đạo hàm để khảo sát và vẽ đồ thị hàm số): Tổng giá trị cực đại và giá trị cực tiểu của hàm số y = x$^3$ - 3x + 12 bằng. Phương án nào sau đây chính xác?\par
\hspace*{1.5em}\textbf{A.} 24\quad\quad \textbf{B.} 0\quad\quad \textbf{C.} 4\quad\quad \textbf{D.} 12\par\vspace{2pt}
\textbf{\color{primary} Câu 30.} Kết hợp các phép biến đổi trong Chương I, hãy tìm kết quả của bài toán: Giá trị lớn nhất của hàm số f(x) = x$^3$ - 3x + 14 trên đoạn [0; 2] bằng:\par
\hspace*{1.5em}\textbf{A.} 16\quad\quad \textbf{B.} 14\quad\quad \textbf{C.} 12\quad\quad \textbf{D.} 18\par\vspace{2pt}
\textbf{\color{primary} Câu 31.} Một bài kiểm tra định kỳ nội dung Chương I nêu câu hỏi: Đường tiệm cận xiên của đồ thị hàm số phân thức y = x + 12 + 3/(x - 1) có phương trình là. Hãy chọn đáp số đúng:\par
\hspace*{1.5em}\textbf{A.} x = 1\quad\quad \textbf{B.} y = x + 12\quad\quad \textbf{C.} y = x + 15\quad\quad \textbf{D.} y = x - 1\par\vspace{2pt}
\textbf{\color{primary} Câu 32.} Tìm lời giải đúng cho bài toán thông hiểu: Tổng giá trị cực đại và giá trị cực tiểu của hàm số y = x$^3$ - 3x + 15 bằng:\par
\hspace*{1.5em}\textbf{A.} 30\quad\quad \textbf{B.} 15\quad\quad \textbf{C.} 4\quad\quad \textbf{D.} 0\par\vspace{2pt}
\textbf{\color{primary} Câu 33.} Sử dụng mối liên hệ giữa các đại lượng trong Chương I: Giá trị lớn nhất của hàm số f(x) = x$^3$ - 3x + 17 trên đoạn [0; 2] bằng:\par
\hspace*{1.5em}\textbf{A.} 15\quad\quad \textbf{B.} 21\quad\quad \textbf{C.} 19\quad\quad \textbf{D.} 17\par\vspace{2pt}
\textbf{\color{primary} Câu 34.} Hãy tính toán và chọn phương án đúng đối với bài tập: Đường tiệm cận xiên của đồ thị hàm số phân thức y = x + 15 + 3/(x - 1) có phương trình là:\par
\hspace*{1.5em}\textbf{A.} x = 1\quad\quad \textbf{B.} y = x - 1\quad\quad \textbf{C.} y = x + 18\quad\quad \textbf{D.} y = x + 15\par\vspace{2pt}
\textbf{\color{primary} Câu 35.} Dựa trên các tính chất trọng tâm của Chương I, xét bài toán: Tổng giá trị cực đại và giá trị cực tiểu của hàm số y = x$^3$ - 3x + 18 bằng:\par
\hspace*{1.5em}\textbf{A.} 36\quad\quad \textbf{B.} 4\quad\quad \textbf{C.} 18\quad\quad \textbf{D.} 0\par\vspace{2pt}
\textbf{\color{primary} Câu 36.} Thực hành giải toán (Ứng dụng đạo hàm để khảo sát và vẽ đồ thị hàm số): Giá trị lớn nhất của hàm số f(x) = x$^3$ - 3x + 20 trên đoạn [0; 2] bằng. Đáp án nào dưới đây thỏa mãn?\par
\hspace*{1.5em}\textbf{A.} 22\quad\quad \textbf{B.} 20\quad\quad \textbf{C.} 24\quad\quad \textbf{D.} 18\par\vspace{2pt}
\textbf{\color{primary} Câu 37.} Kiểm tra kỹ năng biến đổi và suy luận qua bài toán: Đường tiệm cận xiên của đồ thị hàm số phân thức y = x + 18 + 3/(x - 1) có phương trình là:\par
\hspace*{1.5em}\textbf{A.} y = x + 18\quad\quad \textbf{B.} y = x - 1\quad\quad \textbf{C.} y = x + 21\quad\quad \textbf{D.} x = 1\par\vspace{2pt}
\textbf{\color{primary} Câu 38.} Cho giả thiết thuộc phạm vi Chương I: Tổng giá trị cực đại và giá trị cực tiểu của hàm số y = x$^3$ - 3x + 21 bằng. Kết luận đúng là:\par
\hspace*{1.5em}\textbf{A.} 4\quad\quad \textbf{B.} 0\quad\quad \textbf{C.} 21\quad\quad \textbf{D.} 42\par\vspace{2pt}
\textbf{\color{primary} Câu 39.} Rút gọn hoặc tính giá trị theo yêu cầu sau của Chương I: Giá trị lớn nhất của hàm số f(x) = x$^3$ - 3x + 23 trên đoạn [0; 2] bằng:\par
\hspace*{1.5em}\textbf{A.} 21\quad\quad \textbf{B.} 27\quad\quad \textbf{C.} 23\quad\quad \textbf{D.} 25\par\vspace{2pt}
\textbf{\color{primary} Câu 40.} Xác định phương án đúng sau khi thực hiện phép tính / suy luận: Đường tiệm cận xiên của đồ thị hàm số phân thức y = x + 21 + 3/(x - 1) có phương trình là:\par
\hspace*{1.5em}\textbf{A.} y = x + 24\quad\quad \textbf{B.} y = x - 1\quad\quad \textbf{C.} x = 1\quad\quad \textbf{D.} y = x + 21\par\vspace{2pt}
\subsection*{\color{accent} Phần III. Mức độ Khó (Vận dụng \& Vận dụng cao -- Câu 41 đến Câu 50)}
\textbf{\color{primary} Câu 41.} Tìm tất cả các giá trị thực của tham số m để hàm số y = 1/3 $\cdot$ x$^3$ - mx$^2$ + (3m - 2)x + 1 đồng biến trên toàn bộ $\mathbb{R}$:\par
\hspace*{1.5em}\textbf{A.} m $\ge$ 2\quad\quad \textbf{B.} 1 $\le$ m $\le$ 2\quad\quad \textbf{C.} 0 $\le$ m $\le$ 2\quad\quad \textbf{D.} m < 1 hoặc m > 2\par\vspace{2pt}
\textbf{\color{primary} Câu 42.} Bài toán vận dụng (Chương I – Ứng dụng đạo hàm để khảo sát và vẽ đồ thị hàm số): Một cơ sở sản xuất bán x sản phẩm với lợi nhuận P(x) = -x$^2$ + 50x - 100 (triệu đồng). Số sản phẩm x cần bán để lợi nhuận đạt giá trị lớn nhất là. Kết quả đúng là:\par
\hspace*{1.5em}\textbf{A.} x = 100 sản phẩm\quad\quad \textbf{B.} x = 15 sản phẩm\quad\quad \textbf{C.} x = 50 sản phẩm\quad\quad \textbf{D.} x = 25 sản phẩm\par\vspace{2pt}
\textbf{\color{primary} Câu 43.} Vận dụng kiến thức nâng cao trong phạm vi Chương I để giải bài toán: Tìm tất cả các giá trị thực của tham số m để hàm số y = 1/3 $\cdot$ x$^3$ - mx$^2$ + (3m - 2)x + 1 đồng biến trên toàn bộ $\mathbb{R}$:\par
\hspace*{1.5em}\textbf{A.} m < 1 hoặc m > 2\quad\quad \textbf{B.} 1 $\le$ m $\le$ 2\quad\quad \textbf{C.} m $\ge$ 2\quad\quad \textbf{D.} 0 $\le$ m $\le$ 2\par\vspace{2pt}
\textbf{\color{primary} Câu 44.} Phân tích và giải bài toán tổng hợp sau của Chương I: Một cơ sở sản xuất bán x sản phẩm với lợi nhuận P(x) = -x$^2$ + 70x - 100 (triệu đồng). Số sản phẩm x cần bán để lợi nhuận đạt giá trị lớn nhất là:\par
\hspace*{1.5em}\textbf{A.} x = 25 sản phẩm\quad\quad \textbf{B.} x = 35 sản phẩm\quad\quad \textbf{C.} x = 100 sản phẩm\quad\quad \textbf{D.} x = 70 sản phẩm\par\vspace{2pt}
\textbf{\color{primary} Câu 45.} Thử thách tư duy trong phạm vi kiến thức Chương I: Tìm tất cả các giá trị thực của tham số m để hàm số y = 1/3 $\cdot$ x$^3$ - mx$^2$ + (3m - 2)x + 1 đồng biến trên toàn bộ $\mathbb{R}$. Đáp số chính xác là:\par
\hspace*{1.5em}\textbf{A.} m $\ge$ 2\quad\quad \textbf{B.} m < 1 hoặc m > 2\quad\quad \textbf{C.} 1 $\le$ m $\le$ 2\quad\quad \textbf{D.} 0 $\le$ m $\le$ 2\par\vspace{2pt}
\textbf{\color{primary} Câu 46.} Sử dụng phương pháp giải toán của Chương I để tìm đáp án cho bài toán: Một cơ sở sản xuất bán x sản phẩm với lợi nhuận P(x) = -x$^2$ + 50x - 100 (triệu đồng). Số sản phẩm x cần bán để lợi nhuận đạt giá trị lớn nhất là:\par
\hspace*{1.5em}\textbf{A.} x = 15 sản phẩm\quad\quad \textbf{B.} x = 50 sản phẩm\quad\quad \textbf{C.} x = 100 sản phẩm\quad\quad \textbf{D.} x = 25 sản phẩm\par\vspace{2pt}
\textbf{\color{primary} Câu 47.} Xét bài toán vận dụng thực tiễn / nâng cao thuộc Chương I: Tìm tất cả các giá trị thực của tham số m để hàm số y = 1/3 $\cdot$ x$^3$ - mx$^2$ + (3m - 2)x + 1 đồng biến trên toàn bộ $\mathbb{R}$:\par
\hspace*{1.5em}\textbf{A.} m $\ge$ 2\quad\quad \textbf{B.} 1 $\le$ m $\le$ 2\quad\quad \textbf{C.} 0 $\le$ m $\le$ 2\quad\quad \textbf{D.} m < 1 hoặc m > 2\par\vspace{2pt}
\textbf{\color{primary} Câu 48.} Tìm kết quả cuối cùng của bài toán phân loại học sinh khá giỏi (Chương I): Một cơ sở sản xuất bán x sản phẩm với lợi nhuận P(x) = -x$^2$ + 70x - 100 (triệu đồng). Số sản phẩm x cần bán để lợi nhuận đạt giá trị lớn nhất là:\par
\hspace*{1.5em}\textbf{A.} x = 25 sản phẩm\quad\quad \textbf{B.} x = 100 sản phẩm\quad\quad \textbf{C.} x = 70 sản phẩm\quad\quad \textbf{D.} x = 35 sản phẩm\par\vspace{2pt}
\textbf{\color{primary} Câu 49.} Kết hợp các tính chất đã học trong Chương I, hãy giải quyết bài toán: Tìm tất cả các giá trị thực của tham số m để hàm số y = 1/3 $\cdot$ x$^3$ - mx$^2$ + (3m - 2)x + 1 đồng biến trên toàn bộ $\mathbb{R}$:\par
\hspace*{1.5em}\textbf{A.} 1 $\le$ m $\le$ 2\quad\quad \textbf{B.} 0 $\le$ m $\le$ 2\quad\quad \textbf{C.} m $\ge$ 2\quad\quad \textbf{D.} m < 1 hoặc m > 2\par\vspace{2pt}
\textbf{\color{primary} Câu 50.} Câu hỏi vận dụng cao cuối Chương I: Một cơ sở sản xuất bán x sản phẩm với lợi nhuận P(x) = -x$^2$ + 50x - 100 (triệu đồng). Số sản phẩm x cần bán để lợi nhuận đạt giá trị lớn nhất là. Phương án đúng là:\par
\hspace*{1.5em}\textbf{A.} x = 100 sản phẩm\quad\quad \textbf{B.} x = 25 sản phẩm\quad\quad \textbf{C.} x = 15 sản phẩm\quad\quad \textbf{D.} x = 50 sản phẩm\par\vspace{2pt}
\subsection*{\color{primary} Bảng đáp án Chương I: Ứng dụng đạo hàm để khảo sát và vẽ đồ thị hàm số}
\begin{center}
\renewcommand{\arraystretch}{1.25}
\begin{tabularx}{\textwidth}{*{10}{>{\centering\arraybackslash}X}}
\toprule
\textbf{\color{primary} Cột 1} & \textbf{\color{primary} Cột 2} & \textbf{\color{primary} Cột 3} & \textbf{\color{primary} Cột 4} & \textbf{\color{primary} Cột 5} & \textbf{\color{primary} Cột 6} & \textbf{\color{primary} Cột 7} & \textbf{\color{primary} Cột 8} & \textbf{\color{primary} Cột 9} & \textbf{\color{primary} Cột 10} \\
\midrule
\textbf{1.A} & \textbf{2.C} & \textbf{3.C} & \textbf{4.C} & \textbf{5.D} & \textbf{6.B} & \textbf{7.C} & \textbf{8.B} & \textbf{9.B} & \textbf{10.D} \\
\rowcolor{rowtint} \textbf{11.A} & \textbf{12.C} & \textbf{13.A} & \textbf{14.B} & \textbf{15.D} & \textbf{16.A} & \textbf{17.C} & \textbf{18.C} & \textbf{19.C} & \textbf{20.D} \\
\textbf{21.D} & \textbf{22.A} & \textbf{23.C} & \textbf{24.A} & \textbf{25.D} & \textbf{26.A} & \textbf{27.A} & \textbf{28.D} & \textbf{29.A} & \textbf{30.A} \\
\rowcolor{rowtint} \textbf{31.B} & \textbf{32.A} & \textbf{33.C} & \textbf{34.D} & \textbf{35.A} & \textbf{36.A} & \textbf{37.A} & \textbf{38.D} & \textbf{39.D} & \textbf{40.D} \\
\textbf{41.B} & \textbf{42.D} & \textbf{43.B} & \textbf{44.B} & \textbf{45.C} & \textbf{46.D} & \textbf{47.B} & \textbf{48.D} & \textbf{49.A} & \textbf{50.B} \\
\bottomrule
\end{tabularx}
\end{center}\vspace{10pt}
\section*{\color{primary} CHƯƠNG II: VECTƠ VÀ HỆ TRỤC TỌA ĐỘ TRONG KHÔNG GIAN}
\textit{Trọng tâm kiến thức SGK: Vectơ trong không gian (quy tắc hình hộp, trọng tâm), hệ trục tọa độ Oxyz và các phép toán vectơ trong không gian (tọa độ trung điểm, trọng tâm, tích vô hướng, độ dài).}\par\vspace{4pt}
\subsection*{\color{accent} Phần I. Mức độ Dễ (Nhận biết -- Câu 1 đến Câu 20)}
\textbf{\color{primary} Câu 1.} Trong hình hộp ABCD.A'B'C'D', theo quy tắc hình hộp thì tổng ba vectơ $\overrightarrow{AB}$ + $\overrightarrow{AD}$ + $\overrightarrow{AA'}$ bằng:\par
\hspace*{1.5em}\textbf{A.} $\overrightarrow{A'C}$\quad\quad \textbf{B.} $\overrightarrow{AC}$\quad\quad \textbf{C.} $\overrightarrow{AC'}$\quad\quad \textbf{D.} $\overrightarrow{BD'}$\par\vspace{2pt}
\textbf{\color{primary} Câu 2.} Thực hiện yêu cầu nhận biết sau: Trong không gian Oxyz, cho hai điểm A(1; 2; 2) và B(3; -2; 6). Tọa độ trung điểm I của đoạn thẳng AB là.\par
\hspace*{1.5em}\textbf{A.} I(2; 0; 4)\quad\quad \textbf{B.} I(2; 2; 4)\quad\quad \textbf{C.} I(1; -2; 2)\quad\quad \textbf{D.} I(4; 0; 8)\par\vspace{2pt}
\textbf{\color{primary} Câu 3.} Trong bài học thuộc Chương II, xét nội dung: Trong không gian Oxyz, hình chiếu vuông góc của điểm M(4; -3; 5) trên mặt phẳng tọa độ (Oxy) có tọa độ là. Đáp án đúng là:\par
\hspace*{1.5em}\textbf{A.} (4; -3; 0)\quad\quad \textbf{B.} (0; 0; 5)\quad\quad \textbf{C.} (0; -3; 5)\quad\quad \textbf{D.} (4; 0; 5)\par\vspace{2pt}
\textbf{\color{primary} Câu 4.} Dựa vào kiến thức chuẩn SGK của Chương II (Vectơ và hệ trục tọa độ trong không gian), trong hình hộp ABCD.A'B'C'D', theo quy tắc hình hộp thì tổng ba vectơ $\overrightarrow{AB}$ + $\overrightarrow{AD}$ + $\overrightarrow{AA'}$ bằng:\par
\hspace*{1.5em}\textbf{A.} $\overrightarrow{BD'}$\quad\quad \textbf{B.} $\overrightarrow{A'C}$\quad\quad \textbf{C.} $\overrightarrow{AC}$\quad\quad \textbf{D.} $\overrightarrow{AC'}$\par\vspace{2pt}
\textbf{\color{primary} Câu 5.} Chọn kết quả đúng nhất khi xét: Trong không gian Oxyz, cho hai điểm A(1; 2; 5) và B(3; -2; 9). Tọa độ trung điểm I của đoạn thẳng AB là:\par
\hspace*{1.5em}\textbf{A.} I(4; 0; 14)\quad\quad \textbf{B.} I(2; 0; 7)\quad\quad \textbf{C.} I(2; 2; 7)\quad\quad \textbf{D.} I(1; -2; 2)\par\vspace{2pt}
\textbf{\color{primary} Câu 6.} Khi làm bài tập cơ bản về vectơ và hệ trục tọa độ trong không gian, một học sinh gặp câu hỏi: Trong không gian Oxyz, hình chiếu vuông góc của điểm M(7; -3; 5) trên mặt phẳng tọa độ (Oxy) có tọa độ là. Kết quả nào dưới đây đúng?\par
\hspace*{1.5em}\textbf{A.} (0; -3; 5)\quad\quad \textbf{B.} (0; 0; 5)\quad\quad \textbf{C.} (7; 0; 5)\quad\quad \textbf{D.} (7; -3; 0)\par\vspace{2pt}
\textbf{\color{primary} Câu 7.} Xác định giá trị hoặc khẳng định đúng cho bài toán sau: Trong hình hộp ABCD.A'B'C'D', theo quy tắc hình hộp thì tổng ba vectơ $\overrightarrow{AB}$ + $\overrightarrow{AD}$ + $\overrightarrow{AA'}$ bằng:\par
\hspace*{1.5em}\textbf{A.} $\overrightarrow{AC}$\quad\quad \textbf{B.} $\overrightarrow{AC'}$\quad\quad \textbf{C.} $\overrightarrow{A'C}$\quad\quad \textbf{D.} $\overrightarrow{BD'}$\par\vspace{2pt}
\textbf{\color{primary} Câu 8.} Áp dụng trực tiếp định nghĩa và quy tắc trong Chương II: Trong không gian Oxyz, cho hai điểm A(1; 2; 8) và B(3; -2; 12). Tọa độ trung điểm I của đoạn thẳng AB là:\par
\hspace*{1.5em}\textbf{A.} I(4; 0; 20)\quad\quad \textbf{B.} I(1; -2; 2)\quad\quad \textbf{C.} I(2; 2; 10)\quad\quad \textbf{D.} I(2; 0; 10)\par\vspace{2pt}
\textbf{\color{primary} Câu 9.} Kiểm tra nhanh kiến thức nền tảng: Trong không gian Oxyz, hình chiếu vuông góc của điểm M(10; -3; 5) trên mặt phẳng tọa độ (Oxy) có tọa độ là:\par
\hspace*{1.5em}\textbf{A.} (0; 0; 5)\quad\quad \textbf{B.} (10; -3; 0)\quad\quad \textbf{C.} (10; 0; 5)\quad\quad \textbf{D.} (0; -3; 5)\par\vspace{2pt}
\textbf{\color{primary} Câu 10.} Cho dữ kiện nhận biết sau đây: Trong hình hộp ABCD.A'B'C'D', theo quy tắc hình hộp thì tổng ba vectơ $\overrightarrow{AB}$ + $\overrightarrow{AD}$ + $\overrightarrow{AA'}$ bằng. Phương án chính xác là:\par
\hspace*{1.5em}\textbf{A.} $\overrightarrow{AC'}$\quad\quad \textbf{B.} $\overrightarrow{AC}$\quad\quad \textbf{C.} $\overrightarrow{A'C}$\quad\quad \textbf{D.} $\overrightarrow{BD'}$\par\vspace{2pt}
\textbf{\color{primary} Câu 11.} Hãy chỉ ra kết quả đúng trong bốn lựa chọn A, B, C, D đối với yêu cầu: Trong không gian Oxyz, cho hai điểm A(1; 2; 11) và B(3; -2; 15). Tọa độ trung điểm I của đoạn thẳng AB là:\par
\hspace*{1.5em}\textbf{A.} I(4; 0; 26)\quad\quad \textbf{B.} I(2; 2; 13)\quad\quad \textbf{C.} I(1; -2; 2)\quad\quad \textbf{D.} I(2; 0; 13)\par\vspace{2pt}
\textbf{\color{primary} Câu 12.} Câu hỏi ôn tập lý thuyết và kỹ năng cơ bản (Chương II): Trong không gian Oxyz, hình chiếu vuông góc của điểm M(13; -3; 5) trên mặt phẳng tọa độ (Oxy) có tọa độ là:\par
\hspace*{1.5em}\textbf{A.} (0; -3; 5)\quad\quad \textbf{B.} (13; -3; 0)\quad\quad \textbf{C.} (13; 0; 5)\quad\quad \textbf{D.} (0; 0; 5)\par\vspace{2pt}
\textbf{\color{primary} Câu 13.} Theo nội dung sách giáo khoa Kết nối tri thức với cuộc sống (Chương II), trong hình hộp ABCD.A'B'C'D', theo quy tắc hình hộp thì tổng ba vectơ $\overrightarrow{AB}$ + $\overrightarrow{AD}$ + $\overrightarrow{AA'}$ bằng:\par
\hspace*{1.5em}\textbf{A.} $\overrightarrow{AC}$\quad\quad \textbf{B.} $\overrightarrow{BD'}$\quad\quad \textbf{C.} $\overrightarrow{AC'}$\quad\quad \textbf{D.} $\overrightarrow{A'C}$\par\vspace{2pt}
\textbf{\color{primary} Câu 14.} Tìm đáp số đúng của bài toán cơ bản sau: Trong không gian Oxyz, cho hai điểm A(1; 2; 14) và B(3; -2; 18). Tọa độ trung điểm I của đoạn thẳng AB là:\par
\hspace*{1.5em}\textbf{A.} I(4; 0; 32)\quad\quad \textbf{B.} I(2; 2; 16)\quad\quad \textbf{C.} I(2; 0; 16)\quad\quad \textbf{D.} I(1; -2; 2)\par\vspace{2pt}
\textbf{\color{primary} Câu 15.} Xét đối tượng toán học trong Chương II: Trong không gian Oxyz, hình chiếu vuông góc của điểm M(16; -3; 5) trên mặt phẳng tọa độ (Oxy) có tọa độ là. Giá trị hoặc tính chất tương ứng là:\par
\hspace*{1.5em}\textbf{A.} (16; 0; 5)\quad\quad \textbf{B.} (0; -3; 5)\quad\quad \textbf{C.} (16; -3; 0)\quad\quad \textbf{D.} (0; 0; 5)\par\vspace{2pt}
\textbf{\color{primary} Câu 16.} Hoàn thành phép xác định cơ bản sau đây: Trong hình hộp ABCD.A'B'C'D', theo quy tắc hình hộp thì tổng ba vectơ $\overrightarrow{AB}$ + $\overrightarrow{AD}$ + $\overrightarrow{AA'}$ bằng:\par
\hspace*{1.5em}\textbf{A.} $\overrightarrow{BD'}$\quad\quad \textbf{B.} $\overrightarrow{A'C}$\quad\quad \textbf{C.} $\overrightarrow{AC}$\quad\quad \textbf{D.} $\overrightarrow{AC'}$\par\vspace{2pt}
\textbf{\color{primary} Câu 17.} Trong các khẳng định hoặc giá trị dưới đây, đâu là đáp án đúng cho câu hỏi: Trong không gian Oxyz, cho hai điểm A(1; 2; 17) và B(3; -2; 21). Tọa độ trung điểm I của đoạn thẳng AB là?\par
\hspace*{1.5em}\textbf{A.} I(2; 2; 19)\quad\quad \textbf{B.} I(2; 0; 19)\quad\quad \textbf{C.} I(4; 0; 38)\quad\quad \textbf{D.} I(1; -2; 2)\par\vspace{2pt}
\textbf{\color{primary} Câu 18.} Bài tập khởi động Chương II: Trong không gian Oxyz, hình chiếu vuông góc của điểm M(19; -3; 5) trên mặt phẳng tọa độ (Oxy) có tọa độ là:\par
\hspace*{1.5em}\textbf{A.} (0; 0; 5)\quad\quad \textbf{B.} (19; -3; 0)\quad\quad \textbf{C.} (0; -3; 5)\quad\quad \textbf{D.} (19; 0; 5)\par\vspace{2pt}
\textbf{\color{primary} Câu 19.} Quan sát dữ kiện đề bài: Trong hình hộp ABCD.A'B'C'D', theo quy tắc hình hộp thì tổng ba vectơ $\overrightarrow{AB}$ + $\overrightarrow{AD}$ + $\overrightarrow{AA'}$ bằng. Hãy chọn phương án đúng:\par
\hspace*{1.5em}\textbf{A.} $\overrightarrow{AC}$\quad\quad \textbf{B.} $\overrightarrow{BD'}$\quad\quad \textbf{C.} $\overrightarrow{AC'}$\quad\quad \textbf{D.} $\overrightarrow{A'C}$\par\vspace{2pt}
\textbf{\color{primary} Câu 20.} Sử dụng công thức hoặc tính chất đầu bài của Chương II để giải quyết yêu cầu: Trong không gian Oxyz, cho hai điểm A(1; 2; 20) và B(3; -2; 24). Tọa độ trung điểm I của đoạn thẳng AB là:\par
\hspace*{1.5em}\textbf{A.} I(2; 0; 22)\quad\quad \textbf{B.} I(2; 2; 22)\quad\quad \textbf{C.} I(1; -2; 2)\quad\quad \textbf{D.} I(4; 0; 44)\par\vspace{2pt}
\subsection*{\color{accent} Phần II. Mức độ Trung bình (Thông hiểu -- Câu 21 đến Câu 40)}
\textbf{\color{primary} Câu 21.} Trong không gian Oxyz, cho hai vectơ $\vec{u}$ = (2; -1; 1) và $\vec{v}$ = (1; 3; 2). Tích vô hướng $\vec{u}$ $\cdot$ $\vec{v}$ bằng:\par
\hspace*{1.5em}\textbf{A.} -1\quad\quad \textbf{B.} 4\quad\quad \textbf{C.} 1\quad\quad \textbf{D.} 7\par\vspace{2pt}
\textbf{\color{primary} Câu 22.} Vận dụng quy tắc trong Chương II (Vectơ và hệ trục tọa độ trong không gian) để giải bài toán: Trong không gian Oxyz, cho tam giác ABC có A(1; 0; 2), B(2; 3; -1) và C(0; 3; 5). Tọa độ trọng tâm G của tam giác ABC là:\par
\hspace*{1.5em}\textbf{A.} G(3; 6; 6)\quad\quad \textbf{B.} G(1,5; 3; 2)\quad\quad \textbf{C.} G(1; 3; 2)\quad\quad \textbf{D.} G(1; 2; 2)\par\vspace{2pt}
\textbf{\color{primary} Câu 23.} Biến đổi và tính toán theo kiến thức Chương II: Trong không gian Oxyz, độ dài của vectơ $\vec{a}$ = (2; -3; 6) bằng. Kết quả thu được là:\par
\hspace*{1.5em}\textbf{A.} 7\quad\quad \textbf{B.} 49\quad\quad \textbf{C.} $\sqrt{11}$\quad\quad \textbf{D.} 5\par\vspace{2pt}
\textbf{\color{primary} Câu 24.} Xét bài toán mức độ thông hiểu sau: Trong không gian Oxyz, cho hai vectơ $\vec{u}$ = (2; -1; 4) và $\vec{v}$ = (1; 3; 2). Tích vô hướng $\vec{u}$ $\cdot$ $\vec{v}$ bằng. Đáp án đúng là:\par
\hspace*{1.5em}\textbf{A.} 13\quad\quad \textbf{B.} 5\quad\quad \textbf{C.} 7 (1)\quad\quad \textbf{D.} 7\par\vspace{2pt}
\textbf{\color{primary} Câu 25.} Để hoàn thành phiếu học tập Chương II, hãy xác định: Trong không gian Oxyz, cho tam giác ABC có A(1; 0; 2), B(2; 3; -1) và C(0; 3; 14). Tọa độ trọng tâm G của tam giác ABC là:\par
\hspace*{1.5em}\textbf{A.} G(1; 2; 5)\quad\quad \textbf{B.} G(1,5; 3; 5)\quad\quad \textbf{C.} G(1; 3; 5)\quad\quad \textbf{D.} G(3; 6; 15)\par\vspace{2pt}
\textbf{\color{primary} Câu 26.} Phân tích dữ kiện bài toán: Trong không gian Oxyz, độ dài của vectơ $\vec{a}$ = (2; -3; 6) bằng. Giá trị cần tìm bằng:\par
\hspace*{1.5em}\textbf{A.} 5\quad\quad \textbf{B.} 49\quad\quad \textbf{C.} $\sqrt{11}$\quad\quad \textbf{D.} 7\par\vspace{2pt}
\textbf{\color{primary} Câu 27.} Áp dụng tính chất đã học ở Chương II vào bài tập sau: Trong không gian Oxyz, cho hai vectơ $\vec{u}$ = (2; -1; 7) và $\vec{v}$ = (1; 3; 2). Tích vô hướng $\vec{u}$ $\cdot$ $\vec{v}$ bằng:\par
\hspace*{1.5em}\textbf{A.} 13\quad\quad \textbf{B.} 19\quad\quad \textbf{C.} 10\quad\quad \textbf{D.} 11\par\vspace{2pt}
\textbf{\color{primary} Câu 28.} Thực hiện các bước suy luận hợp lý để giải quyết yêu cầu: Trong không gian Oxyz, cho tam giác ABC có A(1; 0; 2), B(2; 3; -1) và C(0; 3; 23). Tọa độ trọng tâm G của tam giác ABC là:\par
\hspace*{1.5em}\textbf{A.} G(1; 2; 8)\quad\quad \textbf{B.} G(3; 6; 24)\quad\quad \textbf{C.} G(1; 3; 8)\quad\quad \textbf{D.} G(1,5; 3; 8)\par\vspace{2pt}
\textbf{\color{primary} Câu 29.} Cho bài toán tự luyện (Vectơ và hệ trục tọa độ trong không gian): Trong không gian Oxyz, độ dài của vectơ $\vec{a}$ = (2; -3; 6) bằng. Phương án nào sau đây chính xác?\par
\hspace*{1.5em}\textbf{A.} 5\quad\quad \textbf{B.} $\sqrt{11}$\quad\quad \textbf{C.} 7\quad\quad \textbf{D.} 49\par\vspace{2pt}
\textbf{\color{primary} Câu 30.} Kết hợp các phép biến đổi trong Chương II, hãy tìm kết quả của bài toán: Trong không gian Oxyz, cho hai vectơ $\vec{u}$ = (2; -1; 10) và $\vec{v}$ = (1; 3; 2). Tích vô hướng $\vec{u}$ $\cdot$ $\vec{v}$ bằng:\par
\hspace*{1.5em}\textbf{A.} 13\quad\quad \textbf{B.} 19\quad\quad \textbf{C.} 17\quad\quad \textbf{D.} 25\par\vspace{2pt}
\textbf{\color{primary} Câu 31.} Một bài kiểm tra định kỳ nội dung Chương II nêu câu hỏi: Trong không gian Oxyz, cho tam giác ABC có A(1; 0; 2), B(2; 3; -1) và C(0; 3; 32). Tọa độ trọng tâm G của tam giác ABC là. Hãy chọn đáp số đúng:\par
\hspace*{1.5em}\textbf{A.} G(1,5; 3; 11)\quad\quad \textbf{B.} G(1; 2; 11)\quad\quad \textbf{C.} G(3; 6; 33)\quad\quad \textbf{D.} G(1; 3; 11)\par\vspace{2pt}
\textbf{\color{primary} Câu 32.} Tìm lời giải đúng cho bài toán thông hiểu: Trong không gian Oxyz, độ dài của vectơ $\vec{a}$ = (2; -3; 6) bằng:\par
\hspace*{1.5em}\textbf{A.} 5\quad\quad \textbf{B.} 7\quad\quad \textbf{C.} $\sqrt{11}$\quad\quad \textbf{D.} 49\par\vspace{2pt}
\textbf{\color{primary} Câu 33.} Sử dụng mối liên hệ giữa các đại lượng trong Chương II: Trong không gian Oxyz, cho hai vectơ $\vec{u}$ = (2; -1; 13) và $\vec{v}$ = (1; 3; 2). Tích vô hướng $\vec{u}$ $\cdot$ $\vec{v}$ bằng:\par
\hspace*{1.5em}\textbf{A.} 23\quad\quad \textbf{B.} 25\quad\quad \textbf{C.} 16\quad\quad \textbf{D.} 31\par\vspace{2pt}
\textbf{\color{primary} Câu 34.} Hãy tính toán và chọn phương án đúng đối với bài tập: Trong không gian Oxyz, cho tam giác ABC có A(1; 0; 2), B(2; 3; -1) và C(0; 3; 41). Tọa độ trọng tâm G của tam giác ABC là:\par
\hspace*{1.5em}\textbf{A.} G(1; 2; 14)\quad\quad \textbf{B.} G(1; 3; 14)\quad\quad \textbf{C.} G(3; 6; 42)\quad\quad \textbf{D.} G(1,5; 3; 14)\par\vspace{2pt}
\textbf{\color{primary} Câu 35.} Dựa trên các tính chất trọng tâm của Chương II, xét bài toán: Trong không gian Oxyz, độ dài của vectơ $\vec{a}$ = (2; -3; 6) bằng:\par
\hspace*{1.5em}\textbf{A.} $\sqrt{11}$\quad\quad \textbf{B.} 49\quad\quad \textbf{C.} 5\quad\quad \textbf{D.} 7\par\vspace{2pt}
\textbf{\color{primary} Câu 36.} Thực hành giải toán (Vectơ và hệ trục tọa độ trong không gian): Trong không gian Oxyz, cho hai vectơ $\vec{u}$ = (2; -1; 16) và $\vec{v}$ = (1; 3; 2). Tích vô hướng $\vec{u}$ $\cdot$ $\vec{v}$ bằng. Đáp án nào dưới đây thỏa mãn?\par
\hspace*{1.5em}\textbf{A.} 29\quad\quad \textbf{B.} 37\quad\quad \textbf{C.} 19\quad\quad \textbf{D.} 31\par\vspace{2pt}
\textbf{\color{primary} Câu 37.} Kiểm tra kỹ năng biến đổi và suy luận qua bài toán: Trong không gian Oxyz, cho tam giác ABC có A(1; 0; 2), B(2; 3; -1) và C(0; 3; 50). Tọa độ trọng tâm G của tam giác ABC là:\par
\hspace*{1.5em}\textbf{A.} G(3; 6; 51)\quad\quad \textbf{B.} G(1,5; 3; 17)\quad\quad \textbf{C.} G(1; 2; 17)\quad\quad \textbf{D.} G(1; 3; 17)\par\vspace{2pt}
\textbf{\color{primary} Câu 38.} Cho giả thiết thuộc phạm vi Chương II: Trong không gian Oxyz, độ dài của vectơ $\vec{a}$ = (2; -3; 6) bằng. Kết luận đúng là:\par
\hspace*{1.5em}\textbf{A.} 7\quad\quad \textbf{B.} 5\quad\quad \textbf{C.} $\sqrt{11}$\quad\quad \textbf{D.} 49\par\vspace{2pt}
\textbf{\color{primary} Câu 39.} Rút gọn hoặc tính giá trị theo yêu cầu sau của Chương II: Trong không gian Oxyz, cho hai vectơ $\vec{u}$ = (2; -1; 19) và $\vec{v}$ = (1; 3; 2). Tích vô hướng $\vec{u}$ $\cdot$ $\vec{v}$ bằng:\par
\hspace*{1.5em}\textbf{A.} 43\quad\quad \textbf{B.} 37\quad\quad \textbf{C.} 22\quad\quad \textbf{D.} 35\par\vspace{2pt}
\textbf{\color{primary} Câu 40.} Xác định phương án đúng sau khi thực hiện phép tính / suy luận: Trong không gian Oxyz, cho tam giác ABC có A(1; 0; 2), B(2; 3; -1) và C(0; 3; 59). Tọa độ trọng tâm G của tam giác ABC là:\par
\hspace*{1.5em}\textbf{A.} G(3; 6; 60)\quad\quad \textbf{B.} G(1; 2; 20)\quad\quad \textbf{C.} G(1,5; 3; 20)\quad\quad \textbf{D.} G(1; 3; 20)\par\vspace{2pt}
\subsection*{\color{accent} Phần III. Mức độ Khó (Vận dụng \& Vận dụng cao -- Câu 41 đến Câu 50)}
\textbf{\color{primary} Câu 41.} Trong không gian Oxyz, cho ba đỉnh của hình bình hành ABCD là A(1; 1; 1), B(2; 3; 4), C(5; 4; 6). Tọa độ đỉnh D là:\par
\hspace*{1.5em}\textbf{A.} D(4; 2; 3)\quad\quad \textbf{B.} D(-2; 0; -3)\quad\quad \textbf{C.} D(6; 6; 9)\quad\quad \textbf{D.} D(4; 2; 5)\par\vspace{2pt}
\textbf{\color{primary} Câu 42.} Bài toán vận dụng (Chương II – Vectơ và hệ trục tọa độ trong không gian): Trong không gian Oxyz, cosin của góc giữa hai vectơ $\vec{u}$ = (1; 2; 2) và $\vec{v}$ = (2; 2; 1) bằng. Kết quả đúng là:\par
\hspace*{1.5em}\textbf{A.} 2/3\quad\quad \textbf{B.} 1/9\quad\quad \textbf{C.} 8/9\quad\quad \textbf{D.} 4/9\par\vspace{2pt}
\textbf{\color{primary} Câu 43.} Vận dụng kiến thức nâng cao trong phạm vi Chương II để giải bài toán: Trong không gian Oxyz, cho ba đỉnh của hình bình hành ABCD là A(1; 1; 1), B(2; 3; 4), C(5; 4; 8). Tọa độ đỉnh D là:\par
\hspace*{1.5em}\textbf{A.} D(4; 2; 7)\quad\quad \textbf{B.} D(-2; 0; -5)\quad\quad \textbf{C.} D(4; 2; 5)\quad\quad \textbf{D.} D(6; 6; 11)\par\vspace{2pt}
\textbf{\color{primary} Câu 44.} Phân tích và giải bài toán tổng hợp sau của Chương II: Trong không gian Oxyz, cosin của góc giữa hai vectơ $\vec{u}$ = (1; 2; 2) và $\vec{v}$ = (2; 2; 1) bằng:\par
\hspace*{1.5em}\textbf{A.} 2/3\quad\quad \textbf{B.} 8/9\quad\quad \textbf{C.} 1/9\quad\quad \textbf{D.} 4/9\par\vspace{2pt}
\textbf{\color{primary} Câu 45.} Thử thách tư duy trong phạm vi kiến thức Chương II: Trong không gian Oxyz, cho ba đỉnh của hình bình hành ABCD là A(1; 1; 1), B(2; 3; 4), C(5; 4; 10). Tọa độ đỉnh D là. Đáp số chính xác là:\par
\hspace*{1.5em}\textbf{A.} D(6; 6; 13)\quad\quad \textbf{B.} D(4; 2; 7)\quad\quad \textbf{C.} D(-2; 0; -7)\quad\quad \textbf{D.} D(4; 2; 9)\par\vspace{2pt}
\textbf{\color{primary} Câu 46.} Sử dụng phương pháp giải toán của Chương II để tìm đáp án cho bài toán: Trong không gian Oxyz, cosin của góc giữa hai vectơ $\vec{u}$ = (1; 2; 2) và $\vec{v}$ = (2; 2; 1) bằng:\par
\hspace*{1.5em}\textbf{A.} 1/9\quad\quad \textbf{B.} 8/9\quad\quad \textbf{C.} 4/9\quad\quad \textbf{D.} 2/3\par\vspace{2pt}
\textbf{\color{primary} Câu 47.} Xét bài toán vận dụng thực tiễn / nâng cao thuộc Chương II: Trong không gian Oxyz, cho ba đỉnh của hình bình hành ABCD là A(1; 1; 1), B(2; 3; 4), C(5; 4; 12). Tọa độ đỉnh D là:\par
\hspace*{1.5em}\textbf{A.} D(4; 2; 11)\quad\quad \textbf{B.} D(-2; 0; -9)\quad\quad \textbf{C.} D(6; 6; 15)\quad\quad \textbf{D.} D(4; 2; 9)\par\vspace{2pt}
\textbf{\color{primary} Câu 48.} Tìm kết quả cuối cùng của bài toán phân loại học sinh khá giỏi (Chương II): Trong không gian Oxyz, cosin của góc giữa hai vectơ $\vec{u}$ = (1; 2; 2) và $\vec{v}$ = (2; 2; 1) bằng:\par
\hspace*{1.5em}\textbf{A.} 8/9\quad\quad \textbf{B.} 4/9\quad\quad \textbf{C.} 1/9\quad\quad \textbf{D.} 2/3\par\vspace{2pt}
\textbf{\color{primary} Câu 49.} Kết hợp các tính chất đã học trong Chương II, hãy giải quyết bài toán: Trong không gian Oxyz, cho ba đỉnh của hình bình hành ABCD là A(1; 1; 1), B(2; 3; 4), C(5; 4; 14). Tọa độ đỉnh D là:\par
\hspace*{1.5em}\textbf{A.} D(4; 2; 13)\quad\quad \textbf{B.} D(4; 2; 11)\quad\quad \textbf{C.} D(6; 6; 17)\quad\quad \textbf{D.} D(-2; 0; -11)\par\vspace{2pt}
\textbf{\color{primary} Câu 50.} Câu hỏi vận dụng cao cuối Chương II: Trong không gian Oxyz, cosin của góc giữa hai vectơ $\vec{u}$ = (1; 2; 2) và $\vec{v}$ = (2; 2; 1) bằng. Phương án đúng là:\par
\hspace*{1.5em}\textbf{A.} 1/9\quad\quad \textbf{B.} 4/9\quad\quad \textbf{C.} 2/3\quad\quad \textbf{D.} 8/9\par\vspace{2pt}
\subsection*{\color{primary} Bảng đáp án Chương II: Vectơ và hệ trục tọa độ trong không gian}
\begin{center}
\renewcommand{\arraystretch}{1.25}
\begin{tabularx}{\textwidth}{*{10}{>{\centering\arraybackslash}X}}
\toprule
\textbf{\color{primary} Cột 1} & \textbf{\color{primary} Cột 2} & \textbf{\color{primary} Cột 3} & \textbf{\color{primary} Cột 4} & \textbf{\color{primary} Cột 5} & \textbf{\color{primary} Cột 6} & \textbf{\color{primary} Cột 7} & \textbf{\color{primary} Cột 8} & \textbf{\color{primary} Cột 9} & \textbf{\color{primary} Cột 10} \\
\midrule
\textbf{1.C} & \textbf{2.A} & \textbf{3.A} & \textbf{4.D} & \textbf{5.B} & \textbf{6.D} & \textbf{7.B} & \textbf{8.D} & \textbf{9.B} & \textbf{10.A} \\
\rowcolor{rowtint} \textbf{11.D} & \textbf{12.B} & \textbf{13.C} & \textbf{14.C} & \textbf{15.C} & \textbf{16.D} & \textbf{17.B} & \textbf{18.B} & \textbf{19.C} & \textbf{20.A} \\
\textbf{21.C} & \textbf{22.D} & \textbf{23.A} & \textbf{24.D} & \textbf{25.A} & \textbf{26.D} & \textbf{27.A} & \textbf{28.A} & \textbf{29.C} & \textbf{30.B} \\
\rowcolor{rowtint} \textbf{31.B} & \textbf{32.B} & \textbf{33.B} & \textbf{34.A} & \textbf{35.D} & \textbf{36.D} & \textbf{37.C} & \textbf{38.A} & \textbf{39.B} & \textbf{40.B} \\
\textbf{41.A} & \textbf{42.C} & \textbf{43.C} & \textbf{44.B} & \textbf{45.B} & \textbf{46.B} & \textbf{47.D} & \textbf{48.A} & \textbf{49.B} & \textbf{50.D} \\
\bottomrule
\end{tabularx}
\end{center}\vspace{10pt}
\section*{\color{primary} CHƯƠNG III: CÁC SỐ ĐẶC TRƯNG ĐO MỨC ĐỘ PHÂN TÁN CỦA MẪU SỐ LIỆU GHÉP NHÓM}
\textit{Trọng tâm kiến thức SGK: Khoảng biến thiên, khoảng tứ phân vị, phương sai và độ lệch chuẩn của mẫu số liệu ghép nhóm và ý nghĩa trong đánh giá độ phân tán/rủi ro.}\par\vspace{4pt}
\subsection*{\color{accent} Phần I. Mức độ Dễ (Nhận biết -- Câu 1 đến Câu 20)}
\textbf{\color{primary} Câu 1.} Cho mẫu số liệu ghép nhóm gồm các nhóm liên tiếp từ [10; 15) đến [35; 40). Khoảng biến thiên R của mẫu số liệu ghép nhóm này bằng:\par
\hspace*{1.5em}\textbf{A.} 5\quad\quad \textbf{B.} 25\quad\quad \textbf{C.} 35\quad\quad \textbf{D.} 30\par\vspace{2pt}
\textbf{\color{primary} Câu 2.} Thực hiện yêu cầu nhận biết sau: Nếu một mẫu số liệu ghép nhóm có tứ phân vị thứ nhất Q$_1$ = 13 và tứ phân vị thứ ba Q$_3$ = 22 thì khoảng tứ phân vị $\Delta_Q$ bằng.\par
\hspace*{1.5em}\textbf{A.} 4,5\quad\quad \textbf{B.} 10\quad\quad \textbf{C.} 9\quad\quad \textbf{D.} 35\par\vspace{2pt}
\textbf{\color{primary} Câu 3.} Trong bài học thuộc Chương III, xét nội dung: Khi so sánh hai mẫu số liệu ghép nhóm có cùng đơn vị đo và số trung bình xấp xỉ nhau, mẫu số liệu nào có độ lệch chuẩn s nhỏ hơn thì. Đáp án đúng là:\par
\begin{itemize}\setlength{\itemsep}{1pt}\setlength{\parskip}{0pt}
  \item[] \textbf{A.} Có mức độ biến động mạnh hơn
  \item[] \textbf{B.} Có khoảng biến thiên lớn hơn
  \item[] \textbf{C.} Có mức độ phân tán thấp hơn (các giá trị tập trung đồng đều hơn quanh số trung bình)
  \item[] \textbf{D.} Có số trung bình lớn hơn
\end{itemize}\vspace{2pt}
\textbf{\color{primary} Câu 4.} Dựa vào kiến thức chuẩn SGK của Chương III (Các số đặc trưng đo mức độ phân tán của mẫu số liệu ghép nhóm), cho mẫu số liệu ghép nhóm gồm các nhóm liên tiếp từ [25; 30) đến [50; 55). Khoảng biến thiên R của mẫu số liệu ghép nhóm này bằng:\par
\hspace*{1.5em}\textbf{A.} 25\quad\quad \textbf{B.} 30\quad\quad \textbf{C.} 35\quad\quad \textbf{D.} 5\par\vspace{2pt}
\textbf{\color{primary} Câu 5.} Chọn kết quả đúng nhất khi xét: Nếu một mẫu số liệu ghép nhóm có tứ phân vị thứ nhất Q$_1$ = 16 và tứ phân vị thứ ba Q$_3$ = 25 thì khoảng tứ phân vị $\Delta_Q$ bằng:\par
\hspace*{1.5em}\textbf{A.} 10\quad\quad \textbf{B.} 4,5\quad\quad \textbf{C.} 41\quad\quad \textbf{D.} 9\par\vspace{2pt}
\textbf{\color{primary} Câu 6.} Khi làm bài tập cơ bản về các số đặc trưng đo mức độ phân tán của mẫu số liệu ghép nhóm, một học sinh gặp câu hỏi: Khi so sánh hai mẫu số liệu ghép nhóm có cùng đơn vị đo và số trung bình xấp xỉ nhau, mẫu số liệu nào có độ lệch chuẩn s nhỏ hơn thì. Kết quả nào dưới đây đúng?\par
\begin{itemize}\setlength{\itemsep}{1pt}\setlength{\parskip}{0pt}
  \item[] \textbf{A.} Có khoảng biến thiên lớn hơn
  \item[] \textbf{B.} Có mức độ phân tán thấp hơn (các giá trị tập trung đồng đều hơn quanh số trung bình)
  \item[] \textbf{C.} Có số trung bình lớn hơn
  \item[] \textbf{D.} Có mức độ biến động mạnh hơn
\end{itemize}\vspace{2pt}
\textbf{\color{primary} Câu 7.} Xác định giá trị hoặc khẳng định đúng cho bài toán sau: Cho mẫu số liệu ghép nhóm gồm các nhóm liên tiếp từ [40; 45) đến [65; 70). Khoảng biến thiên R của mẫu số liệu ghép nhóm này bằng:\par
\hspace*{1.5em}\textbf{A.} 5\quad\quad \textbf{B.} 30\quad\quad \textbf{C.} 25\quad\quad \textbf{D.} 35\par\vspace{2pt}
\textbf{\color{primary} Câu 8.} Áp dụng trực tiếp định nghĩa và quy tắc trong Chương III: Nếu một mẫu số liệu ghép nhóm có tứ phân vị thứ nhất Q$_1$ = 19 và tứ phân vị thứ ba Q$_3$ = 28 thì khoảng tứ phân vị $\Delta_Q$ bằng:\par
\hspace*{1.5em}\textbf{A.} 47\quad\quad \textbf{B.} 9\quad\quad \textbf{C.} 10\quad\quad \textbf{D.} 4,5\par\vspace{2pt}
\textbf{\color{primary} Câu 9.} Kiểm tra nhanh kiến thức nền tảng: Khi so sánh hai mẫu số liệu ghép nhóm có cùng đơn vị đo và số trung bình xấp xỉ nhau, mẫu số liệu nào có độ lệch chuẩn s nhỏ hơn thì:\par
\begin{itemize}\setlength{\itemsep}{1pt}\setlength{\parskip}{0pt}
  \item[] \textbf{A.} Có số trung bình lớn hơn
  \item[] \textbf{B.} Có mức độ biến động mạnh hơn
  \item[] \textbf{C.} Có mức độ phân tán thấp hơn (các giá trị tập trung đồng đều hơn quanh số trung bình)
  \item[] \textbf{D.} Có khoảng biến thiên lớn hơn
\end{itemize}\vspace{2pt}
\textbf{\color{primary} Câu 10.} Cho dữ kiện nhận biết sau đây: Cho mẫu số liệu ghép nhóm gồm các nhóm liên tiếp từ [55; 60) đến [80; 85). Khoảng biến thiên R của mẫu số liệu ghép nhóm này bằng. Phương án chính xác là:\par
\hspace*{1.5em}\textbf{A.} 5\quad\quad \textbf{B.} 25\quad\quad \textbf{C.} 35\quad\quad \textbf{D.} 30\par\vspace{2pt}
\textbf{\color{primary} Câu 11.} Hãy chỉ ra kết quả đúng trong bốn lựa chọn A, B, C, D đối với yêu cầu: Nếu một mẫu số liệu ghép nhóm có tứ phân vị thứ nhất Q$_1$ = 22 và tứ phân vị thứ ba Q$_3$ = 31 thì khoảng tứ phân vị $\Delta_Q$ bằng:\par
\hspace*{1.5em}\textbf{A.} 10\quad\quad \textbf{B.} 53\quad\quad \textbf{C.} 4,5\quad\quad \textbf{D.} 9\par\vspace{2pt}
\textbf{\color{primary} Câu 12.} Câu hỏi ôn tập lý thuyết và kỹ năng cơ bản (Chương III): Khi so sánh hai mẫu số liệu ghép nhóm có cùng đơn vị đo và số trung bình xấp xỉ nhau, mẫu số liệu nào có độ lệch chuẩn s nhỏ hơn thì:\par
\begin{itemize}\setlength{\itemsep}{1pt}\setlength{\parskip}{0pt}
  \item[] \textbf{A.} Có số trung bình lớn hơn
  \item[] \textbf{B.} Có mức độ phân tán thấp hơn (các giá trị tập trung đồng đều hơn quanh số trung bình)
  \item[] \textbf{C.} Có khoảng biến thiên lớn hơn
  \item[] \textbf{D.} Có mức độ biến động mạnh hơn
\end{itemize}\vspace{2pt}
\textbf{\color{primary} Câu 13.} Theo nội dung sách giáo khoa Kết nối tri thức với cuộc sống (Chương III), cho mẫu số liệu ghép nhóm gồm các nhóm liên tiếp từ [70; 75) đến [95; 100). Khoảng biến thiên R của mẫu số liệu ghép nhóm này bằng:\par
\hspace*{1.5em}\textbf{A.} 35\quad\quad \textbf{B.} 25\quad\quad \textbf{C.} 5\quad\quad \textbf{D.} 30\par\vspace{2pt}
\textbf{\color{primary} Câu 14.} Tìm đáp số đúng của bài toán cơ bản sau: Nếu một mẫu số liệu ghép nhóm có tứ phân vị thứ nhất Q$_1$ = 25 và tứ phân vị thứ ba Q$_3$ = 34 thì khoảng tứ phân vị $\Delta_Q$ bằng:\par
\hspace*{1.5em}\textbf{A.} 9\quad\quad \textbf{B.} 10\quad\quad \textbf{C.} 59\quad\quad \textbf{D.} 4,5\par\vspace{2pt}
\textbf{\color{primary} Câu 15.} Xét đối tượng toán học trong Chương III: Khi so sánh hai mẫu số liệu ghép nhóm có cùng đơn vị đo và số trung bình xấp xỉ nhau, mẫu số liệu nào có độ lệch chuẩn s nhỏ hơn thì. Giá trị hoặc tính chất tương ứng là:\par
\begin{itemize}\setlength{\itemsep}{1pt}\setlength{\parskip}{0pt}
  \item[] \textbf{A.} Có số trung bình lớn hơn
  \item[] \textbf{B.} Có mức độ phân tán thấp hơn (các giá trị tập trung đồng đều hơn quanh số trung bình)
  \item[] \textbf{C.} Có mức độ biến động mạnh hơn
  \item[] \textbf{D.} Có khoảng biến thiên lớn hơn
\end{itemize}\vspace{2pt}
\textbf{\color{primary} Câu 16.} Hoàn thành phép xác định cơ bản sau đây: Cho mẫu số liệu ghép nhóm gồm các nhóm liên tiếp từ [85; 90) đến [110; 115). Khoảng biến thiên R của mẫu số liệu ghép nhóm này bằng:\par
\hspace*{1.5em}\textbf{A.} 30\quad\quad \textbf{B.} 25\quad\quad \textbf{C.} 35\quad\quad \textbf{D.} 5\par\vspace{2pt}
\textbf{\color{primary} Câu 17.} Trong các khẳng định hoặc giá trị dưới đây, đâu là đáp án đúng cho câu hỏi: Nếu một mẫu số liệu ghép nhóm có tứ phân vị thứ nhất Q$_1$ = 28 và tứ phân vị thứ ba Q$_3$ = 37 thì khoảng tứ phân vị $\Delta_Q$ bằng?\par
\hspace*{1.5em}\textbf{A.} 65\quad\quad \textbf{B.} 10\quad\quad \textbf{C.} 9\quad\quad \textbf{D.} 4,5\par\vspace{2pt}
\textbf{\color{primary} Câu 18.} Bài tập khởi động Chương III: Khi so sánh hai mẫu số liệu ghép nhóm có cùng đơn vị đo và số trung bình xấp xỉ nhau, mẫu số liệu nào có độ lệch chuẩn s nhỏ hơn thì:\par
\begin{itemize}\setlength{\itemsep}{1pt}\setlength{\parskip}{0pt}
  \item[] \textbf{A.} Có mức độ biến động mạnh hơn
  \item[] \textbf{B.} Có mức độ phân tán thấp hơn (các giá trị tập trung đồng đều hơn quanh số trung bình)
  \item[] \textbf{C.} Có số trung bình lớn hơn
  \item[] \textbf{D.} Có khoảng biến thiên lớn hơn
\end{itemize}\vspace{2pt}
\textbf{\color{primary} Câu 19.} Quan sát dữ kiện đề bài: Cho mẫu số liệu ghép nhóm gồm các nhóm liên tiếp từ [100; 105) đến [125; 130). Khoảng biến thiên R của mẫu số liệu ghép nhóm này bằng. Hãy chọn phương án đúng:\par
\hspace*{1.5em}\textbf{A.} 25\quad\quad \textbf{B.} 35\quad\quad \textbf{C.} 5\quad\quad \textbf{D.} 30\par\vspace{2pt}
\textbf{\color{primary} Câu 20.} Sử dụng công thức hoặc tính chất đầu bài của Chương III để giải quyết yêu cầu: Nếu một mẫu số liệu ghép nhóm có tứ phân vị thứ nhất Q$_1$ = 31 và tứ phân vị thứ ba Q$_3$ = 40 thì khoảng tứ phân vị $\Delta_Q$ bằng:\par
\hspace*{1.5em}\textbf{A.} 9\quad\quad \textbf{B.} 10\quad\quad \textbf{C.} 71\quad\quad \textbf{D.} 4,5\par\vspace{2pt}
\subsection*{\color{accent} Phần II. Mức độ Trung bình (Thông hiểu -- Câu 21 đến Câu 40)}
\textbf{\color{primary} Câu 21.} Một mẫu số liệu ghép nhóm có phương sai s$^2$ = 2.25. Độ lệch chuẩn s của mẫu số liệu ghép nhóm đó bằng:\par
\hspace*{1.5em}\textbf{A.} 3\quad\quad \textbf{B.} 2.25\quad\quad \textbf{C.} 1.5\quad\quad \textbf{D.} 0.75\par\vspace{2pt}
\textbf{\color{primary} Câu 22.} Vận dụng quy tắc trong Chương III (Các số đặc trưng đo mức độ phân tán của mẫu số liệu ghép nhóm) để giải bài toán: Ưu điểm nổi bật của khoảng tứ phân vị $\Delta_Q$ so với khoảng biến thiên R của mẫu số liệu ghép nhóm là gì:\par
\begin{itemize}\setlength{\itemsep}{1pt}\setlength{\parskip}{0pt}
  \item[] \textbf{A.} Tính toán không cần biết tần số của các nhóm
  \item[] \textbf{B.} Luôn bằng phương sai
  \item[] \textbf{C.} Luôn lớn hơn khoảng biến thiên R
  \item[] \textbf{D.} Không bị ảnh hưởng bởi các giá trị ngoại lệ (quá lớn hoặc quá nhỏ ở hai nhóm đầu và cuối)
\end{itemize}\vspace{2pt}
\textbf{\color{primary} Câu 23.} Biến đổi và tính toán theo kiến thức Chương III: Cho mẫu số liệu ghép nhóm gồm 2 nhóm: [0; 10) (giá trị đại diện 5) có tần số 5 và [10; 20) (giá trị đại diện 15) có tần số 5. Phương sai s$^2$ của mẫu số liệu ghép nhóm này bằng. Kết quả thu được là:\par
\hspace*{1.5em}\textbf{A.} 25\quad\quad \textbf{B.} 10\quad\quad \textbf{C.} 5\quad\quad \textbf{D.} 50\par\vspace{2pt}
\textbf{\color{primary} Câu 24.} Xét bài toán mức độ thông hiểu sau: Một mẫu số liệu ghép nhóm có phương sai s$^2$ = 9.00. Độ lệch chuẩn s của mẫu số liệu ghép nhóm đó bằng. Đáp án đúng là:\par
\hspace*{1.5em}\textbf{A.} 1.50\quad\quad \textbf{B.} 9.00\quad\quad \textbf{C.} 3.0\quad\quad \textbf{D.} 6\par\vspace{2pt}
\textbf{\color{primary} Câu 25.} Để hoàn thành phiếu học tập Chương III, hãy xác định: Ưu điểm nổi bật của khoảng tứ phân vị $\Delta_Q$ so với khoảng biến thiên R của mẫu số liệu ghép nhóm là gì:\par
\begin{itemize}\setlength{\itemsep}{1pt}\setlength{\parskip}{0pt}
  \item[] \textbf{A.} Không bị ảnh hưởng bởi các giá trị ngoại lệ (quá lớn hoặc quá nhỏ ở hai nhóm đầu và cuối)
  \item[] \textbf{B.} Tính toán không cần biết tần số của các nhóm
  \item[] \textbf{C.} Luôn bằng phương sai
  \item[] \textbf{D.} Luôn lớn hơn khoảng biến thiên R
\end{itemize}\vspace{2pt}
\textbf{\color{primary} Câu 26.} Phân tích dữ kiện bài toán: Cho mẫu số liệu ghép nhóm gồm 2 nhóm: [0; 10) (giá trị đại diện 5) có tần số 5 và [10; 20) (giá trị đại diện 15) có tần số 5. Phương sai s$^2$ của mẫu số liệu ghép nhóm này bằng. Giá trị cần tìm bằng:\par
\hspace*{1.5em}\textbf{A.} 50\quad\quad \textbf{B.} 25\quad\quad \textbf{C.} 5\quad\quad \textbf{D.} 10\par\vspace{2pt}
\textbf{\color{primary} Câu 27.} Áp dụng tính chất đã học ở Chương III vào bài tập sau: Một mẫu số liệu ghép nhóm có phương sai s$^2$ = 20.25. Độ lệch chuẩn s của mẫu số liệu ghép nhóm đó bằng:\par
\hspace*{1.5em}\textbf{A.} 20.25\quad\quad \textbf{B.} 2.25\quad\quad \textbf{C.} 4.5\quad\quad \textbf{D.} 9\par\vspace{2pt}
\textbf{\color{primary} Câu 28.} Thực hiện các bước suy luận hợp lý để giải quyết yêu cầu: Ưu điểm nổi bật của khoảng tứ phân vị $\Delta_Q$ so với khoảng biến thiên R của mẫu số liệu ghép nhóm là gì:\par
\begin{itemize}\setlength{\itemsep}{1pt}\setlength{\parskip}{0pt}
  \item[] \textbf{A.} Tính toán không cần biết tần số của các nhóm
  \item[] \textbf{B.} Luôn lớn hơn khoảng biến thiên R
  \item[] \textbf{C.} Luôn bằng phương sai
  \item[] \textbf{D.} Không bị ảnh hưởng bởi các giá trị ngoại lệ (quá lớn hoặc quá nhỏ ở hai nhóm đầu và cuối)
\end{itemize}\vspace{2pt}
\textbf{\color{primary} Câu 29.} Cho bài toán tự luyện (Các số đặc trưng đo mức độ phân tán của mẫu số liệu ghép nhóm): Cho mẫu số liệu ghép nhóm gồm 2 nhóm: [0; 10) (giá trị đại diện 5) có tần số 5 và [10; 20) (giá trị đại diện 15) có tần số 5. Phương sai s$^2$ của mẫu số liệu ghép nhóm này bằng. Phương án nào sau đây chính xác?\par
\hspace*{1.5em}\textbf{A.} 5\quad\quad \textbf{B.} 10\quad\quad \textbf{C.} 50\quad\quad \textbf{D.} 25\par\vspace{2pt}
\textbf{\color{primary} Câu 30.} Kết hợp các phép biến đổi trong Chương III, hãy tìm kết quả của bài toán: Một mẫu số liệu ghép nhóm có phương sai s$^2$ = 36.00. Độ lệch chuẩn s của mẫu số liệu ghép nhóm đó bằng:\par
\hspace*{1.5em}\textbf{A.} 36.00\quad\quad \textbf{B.} 6.0\quad\quad \textbf{C.} 3.00\quad\quad \textbf{D.} 12\par\vspace{2pt}
\textbf{\color{primary} Câu 31.} Một bài kiểm tra định kỳ nội dung Chương III nêu câu hỏi: Ưu điểm nổi bật của khoảng tứ phân vị $\Delta_Q$ so với khoảng biến thiên R của mẫu số liệu ghép nhóm là gì. Hãy chọn đáp số đúng:\par
\begin{itemize}\setlength{\itemsep}{1pt}\setlength{\parskip}{0pt}
  \item[] \textbf{A.} Luôn lớn hơn khoảng biến thiên R
  \item[] \textbf{B.} Tính toán không cần biết tần số của các nhóm
  \item[] \textbf{C.} Luôn bằng phương sai
  \item[] \textbf{D.} Không bị ảnh hưởng bởi các giá trị ngoại lệ (quá lớn hoặc quá nhỏ ở hai nhóm đầu và cuối)
\end{itemize}\vspace{2pt}
\textbf{\color{primary} Câu 32.} Tìm lời giải đúng cho bài toán thông hiểu: Cho mẫu số liệu ghép nhóm gồm 2 nhóm: [0; 10) (giá trị đại diện 5) có tần số 5 và [10; 20) (giá trị đại diện 15) có tần số 5. Phương sai s$^2$ của mẫu số liệu ghép nhóm này bằng:\par
\hspace*{1.5em}\textbf{A.} 5\quad\quad \textbf{B.} 10\quad\quad \textbf{C.} 25\quad\quad \textbf{D.} 50\par\vspace{2pt}
\textbf{\color{primary} Câu 33.} Sử dụng mối liên hệ giữa các đại lượng trong Chương III: Một mẫu số liệu ghép nhóm có phương sai s$^2$ = 56.25. Độ lệch chuẩn s của mẫu số liệu ghép nhóm đó bằng:\par
\hspace*{1.5em}\textbf{A.} 15\quad\quad \textbf{B.} 56.25\quad\quad \textbf{C.} 7.5\quad\quad \textbf{D.} 3.75\par\vspace{2pt}
\textbf{\color{primary} Câu 34.} Hãy tính toán và chọn phương án đúng đối với bài tập: Ưu điểm nổi bật của khoảng tứ phân vị $\Delta_Q$ so với khoảng biến thiên R của mẫu số liệu ghép nhóm là gì:\par
\begin{itemize}\setlength{\itemsep}{1pt}\setlength{\parskip}{0pt}
  \item[] \textbf{A.} Luôn lớn hơn khoảng biến thiên R
  \item[] \textbf{B.} Không bị ảnh hưởng bởi các giá trị ngoại lệ (quá lớn hoặc quá nhỏ ở hai nhóm đầu và cuối)
  \item[] \textbf{C.} Luôn bằng phương sai
  \item[] \textbf{D.} Tính toán không cần biết tần số của các nhóm
\end{itemize}\vspace{2pt}
\textbf{\color{primary} Câu 35.} Dựa trên các tính chất trọng tâm của Chương III, xét bài toán: Cho mẫu số liệu ghép nhóm gồm 2 nhóm: [0; 10) (giá trị đại diện 5) có tần số 5 và [10; 20) (giá trị đại diện 15) có tần số 5. Phương sai s$^2$ của mẫu số liệu ghép nhóm này bằng:\par
\hspace*{1.5em}\textbf{A.} 10\quad\quad \textbf{B.} 5\quad\quad \textbf{C.} 50\quad\quad \textbf{D.} 25\par\vspace{2pt}
\textbf{\color{primary} Câu 36.} Thực hành giải toán (Các số đặc trưng đo mức độ phân tán của mẫu số liệu ghép nhóm): Một mẫu số liệu ghép nhóm có phương sai s$^2$ = 81.00. Độ lệch chuẩn s của mẫu số liệu ghép nhóm đó bằng. Đáp án nào dưới đây thỏa mãn?\par
\hspace*{1.5em}\textbf{A.} 81.00\quad\quad \textbf{B.} 18\quad\quad \textbf{C.} 9.0\quad\quad \textbf{D.} 4.50\par\vspace{2pt}
\textbf{\color{primary} Câu 37.} Kiểm tra kỹ năng biến đổi và suy luận qua bài toán: Ưu điểm nổi bật của khoảng tứ phân vị $\Delta_Q$ so với khoảng biến thiên R của mẫu số liệu ghép nhóm là gì:\par
\begin{itemize}\setlength{\itemsep}{1pt}\setlength{\parskip}{0pt}
  \item[] \textbf{A.} Luôn bằng phương sai
  \item[] \textbf{B.} Không bị ảnh hưởng bởi các giá trị ngoại lệ (quá lớn hoặc quá nhỏ ở hai nhóm đầu và cuối)
  \item[] \textbf{C.} Tính toán không cần biết tần số của các nhóm
  \item[] \textbf{D.} Luôn lớn hơn khoảng biến thiên R
\end{itemize}\vspace{2pt}
\textbf{\color{primary} Câu 38.} Cho giả thiết thuộc phạm vi Chương III: Cho mẫu số liệu ghép nhóm gồm 2 nhóm: [0; 10) (giá trị đại diện 5) có tần số 5 và [10; 20) (giá trị đại diện 15) có tần số 5. Phương sai s$^2$ của mẫu số liệu ghép nhóm này bằng. Kết luận đúng là:\par
\hspace*{1.5em}\textbf{A.} 50\quad\quad \textbf{B.} 5\quad\quad \textbf{C.} 25\quad\quad \textbf{D.} 10\par\vspace{2pt}
\textbf{\color{primary} Câu 39.} Rút gọn hoặc tính giá trị theo yêu cầu sau của Chương III: Một mẫu số liệu ghép nhóm có phương sai s$^2$ = 110.25. Độ lệch chuẩn s của mẫu số liệu ghép nhóm đó bằng:\par
\hspace*{1.5em}\textbf{A.} 5.25\quad\quad \textbf{B.} 110.25\quad\quad \textbf{C.} 21\quad\quad \textbf{D.} 10.5\par\vspace{2pt}
\textbf{\color{primary} Câu 40.} Xác định phương án đúng sau khi thực hiện phép tính / suy luận: Ưu điểm nổi bật của khoảng tứ phân vị $\Delta_Q$ so với khoảng biến thiên R của mẫu số liệu ghép nhóm là gì:\par
\begin{itemize}\setlength{\itemsep}{1pt}\setlength{\parskip}{0pt}
  \item[] \textbf{A.} Không bị ảnh hưởng bởi các giá trị ngoại lệ (quá lớn hoặc quá nhỏ ở hai nhóm đầu và cuối)
  \item[] \textbf{B.} Luôn lớn hơn khoảng biến thiên R
  \item[] \textbf{C.} Tính toán không cần biết tần số của các nhóm
  \item[] \textbf{D.} Luôn bằng phương sai
\end{itemize}\vspace{2pt}
\subsection*{\color{accent} Phần III. Mức độ Khó (Vận dụng \& Vận dụng cao -- Câu 41 đến Câu 50)}
\textbf{\color{primary} Câu 41.} Cho mẫu số liệu ghép nhóm gồm 3 nhóm: [0; 10) có tần số 2; [10; 20) có tần số 6; [20; 30) có tần số 2 (giá trị đại diện lần lượt là 5; 15; 25). Phương sai s$^2$ của mẫu số liệu ghép nhóm bằng:\par
\hspace*{1.5em}\textbf{A.} 40\quad\quad \textbf{B.} 20\quad\quad \textbf{C.} 100\quad\quad \textbf{D.} $\sqrt{40}$\par\vspace{2pt}
\textbf{\color{primary} Câu 42.} Bài toán vận dụng (Chương III – Các số đặc trưng đo mức độ phân tán của mẫu số liệu ghép nhóm): Nếu cộng thêm cùng một hằng số c $\ne$ 0 vào mọi đầu mút của các nhóm trong mẫu số liệu ghép nhóm (giữ nguyên tần số các nhóm) thì phương sai s$^2$ và độ lệch chuẩn s thay đổi như thế nào. Kết quả đúng là:\par
\begin{itemize}\setlength{\itemsep}{1pt}\setlength{\parskip}{0pt}
  \item[] \textbf{A.} Tăng thêm c đơn vị
  \item[] \textbf{B.} Nhân lên c lần
  \item[] \textbf{C.} Tăng thêm c$^2$ đơn vị
  \item[] \textbf{D.} Không thay đổi (giữ nguyên giá trị ban đầu)
\end{itemize}\vspace{2pt}
\textbf{\color{primary} Câu 43.} Vận dụng kiến thức nâng cao trong phạm vi Chương III để giải bài toán: Cho mẫu số liệu ghép nhóm gồm 3 nhóm: [20; 30) tần số 2; [30; 40) tần số 6; [40; 50) tần số 2. Phương sai s$^2$ của mẫu số liệu ghép nhóm này bằng:\par
\hspace*{1.5em}\textbf{A.} 100\quad\quad \textbf{B.} 60\quad\quad \textbf{C.} 40\quad\quad \textbf{D.} 20\par\vspace{2pt}
\textbf{\color{primary} Câu 44.} Phân tích và giải bài toán tổng hợp sau của Chương III: Cho mẫu số liệu ghép nhóm gồm 3 nhóm: [0; 10) có tần số 2; [10; 20) có tần số 6; [20; 30) có tần số 2 (giá trị đại diện lần lượt là 5; 15; 25). Phương sai s$^2$ của mẫu số liệu ghép nhóm bằng:\par
\hspace*{1.5em}\textbf{A.} 20\quad\quad \textbf{B.} 40\quad\quad \textbf{C.} 100\quad\quad \textbf{D.} $\sqrt{40}$\par\vspace{2pt}
\textbf{\color{primary} Câu 45.} Thử thách tư duy trong phạm vi kiến thức Chương III: Nếu cộng thêm cùng một hằng số c $\ne$ 0 vào mọi đầu mút của các nhóm trong mẫu số liệu ghép nhóm (giữ nguyên tần số các nhóm) thì phương sai s$^2$ và độ lệch chuẩn s thay đổi như thế nào. Đáp số chính xác là:\par
\begin{itemize}\setlength{\itemsep}{1pt}\setlength{\parskip}{0pt}
  \item[] \textbf{A.} Tăng thêm c đơn vị
  \item[] \textbf{B.} Nhân lên c lần
  \item[] \textbf{C.} Không thay đổi (giữ nguyên giá trị ban đầu)
  \item[] \textbf{D.} Tăng thêm c$^2$ đơn vị
\end{itemize}\vspace{2pt}
\textbf{\color{primary} Câu 46.} Sử dụng phương pháp giải toán của Chương III để tìm đáp án cho bài toán: Cho mẫu số liệu ghép nhóm gồm 3 nhóm: [50; 60) tần số 2; [60; 70) tần số 6; [70; 80) tần số 2. Phương sai s$^2$ của mẫu số liệu ghép nhóm này bằng:\par
\hspace*{1.5em}\textbf{A.} 20\quad\quad \textbf{B.} 90\quad\quad \textbf{C.} 40\quad\quad \textbf{D.} 100\par\vspace{2pt}
\textbf{\color{primary} Câu 47.} Xét bài toán vận dụng thực tiễn / nâng cao thuộc Chương III: Cho mẫu số liệu ghép nhóm gồm 3 nhóm: [0; 10) có tần số 2; [10; 20) có tần số 6; [20; 30) có tần số 2 (giá trị đại diện lần lượt là 5; 15; 25). Phương sai s$^2$ của mẫu số liệu ghép nhóm bằng:\par
\hspace*{1.5em}\textbf{A.} $\sqrt{40}$\quad\quad \textbf{B.} 40\quad\quad \textbf{C.} 100\quad\quad \textbf{D.} 20\par\vspace{2pt}
\textbf{\color{primary} Câu 48.} Tìm kết quả cuối cùng của bài toán phân loại học sinh khá giỏi (Chương III): Nếu cộng thêm cùng một hằng số c $\ne$ 0 vào mọi đầu mút của các nhóm trong mẫu số liệu ghép nhóm (giữ nguyên tần số các nhóm) thì phương sai s$^2$ và độ lệch chuẩn s thay đổi như thế nào:\par
\begin{itemize}\setlength{\itemsep}{1pt}\setlength{\parskip}{0pt}
  \item[] \textbf{A.} Tăng thêm c đơn vị
  \item[] \textbf{B.} Nhân lên c lần
  \item[] \textbf{C.} Không thay đổi (giữ nguyên giá trị ban đầu)
  \item[] \textbf{D.} Tăng thêm c$^2$ đơn vị
\end{itemize}\vspace{2pt}
\textbf{\color{primary} Câu 49.} Kết hợp các tính chất đã học trong Chương III, hãy giải quyết bài toán: Cho mẫu số liệu ghép nhóm gồm 3 nhóm: [80; 90) tần số 2; [90; 100) tần số 6; [100; 110) tần số 2. Phương sai s$^2$ của mẫu số liệu ghép nhóm này bằng:\par
\hspace*{1.5em}\textbf{A.} 100\quad\quad \textbf{B.} 40\quad\quad \textbf{C.} 120\quad\quad \textbf{D.} 20\par\vspace{2pt}
\textbf{\color{primary} Câu 50.} Câu hỏi vận dụng cao cuối Chương III: Cho mẫu số liệu ghép nhóm gồm 3 nhóm: [0; 10) có tần số 2; [10; 20) có tần số 6; [20; 30) có tần số 2 (giá trị đại diện lần lượt là 5; 15; 25). Phương sai s$^2$ của mẫu số liệu ghép nhóm bằng. Phương án đúng là:\par
\hspace*{1.5em}\textbf{A.} $\sqrt{40}$\quad\quad \textbf{B.} 20\quad\quad \textbf{C.} 100\quad\quad \textbf{D.} 40\par\vspace{2pt}
\subsection*{\color{primary} Bảng đáp án Chương III: Các số đặc trưng đo mức độ phân tán của mẫu số liệu ghép nhóm}
\begin{center}
\renewcommand{\arraystretch}{1.25}
\begin{tabularx}{\textwidth}{*{10}{>{\centering\arraybackslash}X}}
\toprule
\textbf{\color{primary} Cột 1} & \textbf{\color{primary} Cột 2} & \textbf{\color{primary} Cột 3} & \textbf{\color{primary} Cột 4} & \textbf{\color{primary} Cột 5} & \textbf{\color{primary} Cột 6} & \textbf{\color{primary} Cột 7} & \textbf{\color{primary} Cột 8} & \textbf{\color{primary} Cột 9} & \textbf{\color{primary} Cột 10} \\
\midrule
\textbf{1.D} & \textbf{2.C} & \textbf{3.C} & \textbf{4.B} & \textbf{5.D} & \textbf{6.B} & \textbf{7.B} & \textbf{8.B} & \textbf{9.C} & \textbf{10.D} \\
\rowcolor{rowtint} \textbf{11.D} & \textbf{12.B} & \textbf{13.D} & \textbf{14.A} & \textbf{15.B} & \textbf{16.A} & \textbf{17.C} & \textbf{18.B} & \textbf{19.D} & \textbf{20.A} \\
\textbf{21.C} & \textbf{22.D} & \textbf{23.A} & \textbf{24.C} & \textbf{25.A} & \textbf{26.B} & \textbf{27.C} & \textbf{28.D} & \textbf{29.D} & \textbf{30.B} \\
\rowcolor{rowtint} \textbf{31.D} & \textbf{32.C} & \textbf{33.C} & \textbf{34.B} & \textbf{35.D} & \textbf{36.C} & \textbf{37.B} & \textbf{38.C} & \textbf{39.D} & \textbf{40.A} \\
\textbf{41.A} & \textbf{42.D} & \textbf{43.C} & \textbf{44.B} & \textbf{45.C} & \textbf{46.C} & \textbf{47.B} & \textbf{48.C} & \textbf{49.B} & \textbf{50.D} \\
\bottomrule
\end{tabularx}
\end{center}\vspace{10pt}
\section*{\color{primary} CHƯƠNG IV: NGUYÊN HÀM VÀ TÍCH PHÂN}
\textit{Trọng tâm kiến thức SGK: Nguyên hàm của hàm số sơ cấp, tích phân xác định và tính chất, ứng dụng hình học của tích phân để tính diện tích hình phẳng và thể tích vật thể tròn xoay.}\par\vspace{4pt}
\subsection*{\color{accent} Phần I. Mức độ Dễ (Nhận biết -- Câu 1 đến Câu 20)}
\textbf{\color{primary} Câu 1.} Họ tất cả các nguyên hàm của hàm số f(x) = 3x$^2$ + 2x là:\par
\hspace*{1.5em}\textbf{A.} 3x$^3$ + 2x$^2$ + C\quad\quad \textbf{B.} 6x + 2 + C\quad\quad \textbf{C.} x$^3$ + 1x$^2$ + C\quad\quad \textbf{D.} x$^3$ + 2x$^2$ + C\par\vspace{2pt}
\textbf{\color{primary} Câu 2.} Thực hiện yêu cầu nhận biết sau: Giá trị của tích phân xác định $\int_{0}^{2}$ (2x + 2) dx bằng.\par
\hspace*{1.5em}\textbf{A.} 6\quad\quad \textbf{B.} 8\quad\quad \textbf{C.} 4\quad\quad \textbf{D.} 8 (1)\par\vspace{2pt}
\textbf{\color{primary} Câu 3.} Trong bài học thuộc Chương IV, xét nội dung: Công thức tính diện tích S của hình phẳng giới hạn bởi đồ thị hàm số liên tục y = f(x), trục hoành Ox và hai đường thẳng x = a, x = b (a < b) là. Đáp án đúng là:\par
\hspace*{1.5em}\textbf{A.} S = $\int_{a}^{b}$ |f(x)| dx\quad\quad \textbf{B.} S = $\int_{a}^{b}$ f(x) dx\quad\quad \textbf{C.} S = $\int_{a}^{b}$ [f(x)]$^2$ dx\quad\quad \textbf{D.} S = $\pi$ $\int_{a}^{b}$ [f(x)]$^2$ dx\par\vspace{2pt}
\textbf{\color{primary} Câu 4.} Dựa vào kiến thức chuẩn SGK của Chương IV (Nguyên hàm và tích phân), họ tất cả các nguyên hàm của hàm số f(x) = 3x$^2$ + 8x là:\par
\hspace*{1.5em}\textbf{A.} x$^3$ + 4x$^2$ + C\quad\quad \textbf{B.} 6x + 8 + C\quad\quad \textbf{C.} x$^3$ + 8x$^2$ + C\quad\quad \textbf{D.} 3x$^3$ + 8x$^2$ + C\par\vspace{2pt}
\textbf{\color{primary} Câu 5.} Chọn kết quả đúng nhất khi xét: Giá trị của tích phân xác định $\int_{0}^{2}$ (2x + 5) dx bằng:\par
\hspace*{1.5em}\textbf{A.} 14\quad\quad \textbf{B.} 4\quad\quad \textbf{C.} 12\quad\quad \textbf{D.} 9\par\vspace{2pt}
\textbf{\color{primary} Câu 6.} Khi làm bài tập cơ bản về nguyên hàm và tích phân, một học sinh gặp câu hỏi: Công thức tính diện tích S của hình phẳng giới hạn bởi đồ thị hàm số liên tục y = f(x), trục hoành Ox và hai đường thẳng x = a, x = b (a < b) là. Kết quả nào dưới đây đúng?\par
\hspace*{1.5em}\textbf{A.} S = $\int_{a}^{b}$ [f(x)]$^2$ dx\quad\quad \textbf{B.} S = $\pi$ $\int_{a}^{b}$ [f(x)]$^2$ dx\quad\quad \textbf{C.} S = $\int_{a}^{b}$ |f(x)| dx\quad\quad \textbf{D.} S = $\int_{a}^{b}$ f(x) dx\par\vspace{2pt}
\textbf{\color{primary} Câu 7.} Xác định giá trị hoặc khẳng định đúng cho bài toán sau: Họ tất cả các nguyên hàm của hàm số f(x) = 3x$^2$ + 14x là:\par
\hspace*{1.5em}\textbf{A.} 3x$^3$ + 14x$^2$ + C\quad\quad \textbf{B.} x$^3$ + 7x$^2$ + C\quad\quad \textbf{C.} 6x + 14 + C\quad\quad \textbf{D.} x$^3$ + 14x$^2$ + C\par\vspace{2pt}
\textbf{\color{primary} Câu 8.} Áp dụng trực tiếp định nghĩa và quy tắc trong Chương IV: Giá trị của tích phân xác định $\int_{0}^{2}$ (2x + 8) dx bằng:\par
\hspace*{1.5em}\textbf{A.} 18\quad\quad \textbf{B.} 12\quad\quad \textbf{C.} 4\quad\quad \textbf{D.} 20\par\vspace{2pt}
\textbf{\color{primary} Câu 9.} Kiểm tra nhanh kiến thức nền tảng: Công thức tính diện tích S của hình phẳng giới hạn bởi đồ thị hàm số liên tục y = f(x), trục hoành Ox và hai đường thẳng x = a, x = b (a < b) là:\par
\hspace*{1.5em}\textbf{A.} S = $\pi$ $\int_{a}^{b}$ [f(x)]$^2$ dx\quad\quad \textbf{B.} S = $\int_{a}^{b}$ [f(x)]$^2$ dx\quad\quad \textbf{C.} S = $\int_{a}^{b}$ |f(x)| dx\quad\quad \textbf{D.} S = $\int_{a}^{b}$ f(x) dx\par\vspace{2pt}
\textbf{\color{primary} Câu 10.} Cho dữ kiện nhận biết sau đây: Họ tất cả các nguyên hàm của hàm số f(x) = 3x$^2$ + 20x là. Phương án chính xác là:\par
\hspace*{1.5em}\textbf{A.} 6x + 20 + C\quad\quad \textbf{B.} 3x$^3$ + 20x$^2$ + C\quad\quad \textbf{C.} x$^3$ + 20x$^2$ + C\quad\quad \textbf{D.} x$^3$ + 10x$^2$ + C\par\vspace{2pt}
\textbf{\color{primary} Câu 11.} Hãy chỉ ra kết quả đúng trong bốn lựa chọn A, B, C, D đối với yêu cầu: Giá trị của tích phân xác định $\int_{0}^{2}$ (2x + 11) dx bằng:\par
\hspace*{1.5em}\textbf{A.} 24\quad\quad \textbf{B.} 4\quad\quad \textbf{C.} 26\quad\quad \textbf{D.} 15\par\vspace{2pt}
\textbf{\color{primary} Câu 12.} Câu hỏi ôn tập lý thuyết và kỹ năng cơ bản (Chương IV): Công thức tính diện tích S của hình phẳng giới hạn bởi đồ thị hàm số liên tục y = f(x), trục hoành Ox và hai đường thẳng x = a, x = b (a < b) là:\par
\hspace*{1.5em}\textbf{A.} S = $\int_{a}^{b}$ |f(x)| dx\quad\quad \textbf{B.} S = $\int_{a}^{b}$ [f(x)]$^2$ dx\quad\quad \textbf{C.} S = $\int_{a}^{b}$ f(x) dx\quad\quad \textbf{D.} S = $\pi$ $\int_{a}^{b}$ [f(x)]$^2$ dx\par\vspace{2pt}
\textbf{\color{primary} Câu 13.} Theo nội dung sách giáo khoa Kết nối tri thức với cuộc sống (Chương IV), họ tất cả các nguyên hàm của hàm số f(x) = 3x$^2$ + 26x là:\par
\hspace*{1.5em}\textbf{A.} 3x$^3$ + 26x$^2$ + C\quad\quad \textbf{B.} x$^3$ + 13x$^2$ + C\quad\quad \textbf{C.} x$^3$ + 26x$^2$ + C\quad\quad \textbf{D.} 6x + 26 + C\par\vspace{2pt}
\textbf{\color{primary} Câu 14.} Tìm đáp số đúng của bài toán cơ bản sau: Giá trị của tích phân xác định $\int_{0}^{2}$ (2x + 14) dx bằng:\par
\hspace*{1.5em}\textbf{A.} 30\quad\quad \textbf{B.} 4\quad\quad \textbf{C.} 18\quad\quad \textbf{D.} 32\par\vspace{2pt}
\textbf{\color{primary} Câu 15.} Xét đối tượng toán học trong Chương IV: Công thức tính diện tích S của hình phẳng giới hạn bởi đồ thị hàm số liên tục y = f(x), trục hoành Ox và hai đường thẳng x = a, x = b (a < b) là. Giá trị hoặc tính chất tương ứng là:\par
\hspace*{1.5em}\textbf{A.} S = $\int_{a}^{b}$ |f(x)| dx\quad\quad \textbf{B.} S = $\pi$ $\int_{a}^{b}$ [f(x)]$^2$ dx\quad\quad \textbf{C.} S = $\int_{a}^{b}$ [f(x)]$^2$ dx\quad\quad \textbf{D.} S = $\int_{a}^{b}$ f(x) dx\par\vspace{2pt}
\textbf{\color{primary} Câu 16.} Hoàn thành phép xác định cơ bản sau đây: Họ tất cả các nguyên hàm của hàm số f(x) = 3x$^2$ + 32x là:\par
\hspace*{1.5em}\textbf{A.} 3x$^3$ + 32x$^2$ + C\quad\quad \textbf{B.} x$^3$ + 16x$^2$ + C\quad\quad \textbf{C.} 6x + 32 + C\quad\quad \textbf{D.} x$^3$ + 32x$^2$ + C\par\vspace{2pt}
\textbf{\color{primary} Câu 17.} Trong các khẳng định hoặc giá trị dưới đây, đâu là đáp án đúng cho câu hỏi: Giá trị của tích phân xác định $\int_{0}^{2}$ (2x + 17) dx bằng?\par
\hspace*{1.5em}\textbf{A.} 38\quad\quad \textbf{B.} 4\quad\quad \textbf{C.} 36\quad\quad \textbf{D.} 21\par\vspace{2pt}
\textbf{\color{primary} Câu 18.} Bài tập khởi động Chương IV: Công thức tính diện tích S của hình phẳng giới hạn bởi đồ thị hàm số liên tục y = f(x), trục hoành Ox và hai đường thẳng x = a, x = b (a < b) là:\par
\hspace*{1.5em}\textbf{A.} S = $\int_{a}^{b}$ [f(x)]$^2$ dx\quad\quad \textbf{B.} S = $\int_{a}^{b}$ |f(x)| dx\quad\quad \textbf{C.} S = $\pi$ $\int_{a}^{b}$ [f(x)]$^2$ dx\quad\quad \textbf{D.} S = $\int_{a}^{b}$ f(x) dx\par\vspace{2pt}
\textbf{\color{primary} Câu 19.} Quan sát dữ kiện đề bài: Họ tất cả các nguyên hàm của hàm số f(x) = 3x$^2$ + 38x là. Hãy chọn phương án đúng:\par
\hspace*{1.5em}\textbf{A.} x$^3$ + 38x$^2$ + C\quad\quad \textbf{B.} x$^3$ + 19x$^2$ + C\quad\quad \textbf{C.} 6x + 38 + C\quad\quad \textbf{D.} 3x$^3$ + 38x$^2$ + C\par\vspace{2pt}
\textbf{\color{primary} Câu 20.} Sử dụng công thức hoặc tính chất đầu bài của Chương IV để giải quyết yêu cầu: Giá trị của tích phân xác định $\int_{0}^{2}$ (2x + 20) dx bằng:\par
\hspace*{1.5em}\textbf{A.} 42\quad\quad \textbf{B.} 44\quad\quad \textbf{C.} 4\quad\quad \textbf{D.} 24\par\vspace{2pt}
\subsection*{\color{accent} Phần II. Mức độ Trung bình (Thông hiểu -- Câu 21 đến Câu 40)}
\textbf{\color{primary} Câu 21.} Cho $\int_{1}^{3}$ f(x) dx = 4 và $\int_{1}^{3}$ g(x) dx = 2. Giá trị của tích phân I = $\int_{1}^{3}$ [2f(x) - 3g(x)] dx bằng:\par
\hspace*{1.5em}\textbf{A.} 2\quad\quad \textbf{B.} 6\quad\quad \textbf{C.} 14\quad\quad \textbf{D.} 2 (1)\par\vspace{2pt}
\textbf{\color{primary} Câu 22.} Vận dụng quy tắc trong Chương IV (Nguyên hàm và tích phân) để giải bài toán: Tìm nguyên hàm F(x) của hàm số f(x) = 2x + 3 biết F(0) = 5:\par
\hspace*{1.5em}\textbf{A.} F(x) = x$^2$ + 5\quad\quad \textbf{B.} F(x) = x$^2$ + 3x - 5\quad\quad \textbf{C.} F(x) = 2x$^2$ + 3x + 5\quad\quad \textbf{D.} F(x) = x$^2$ + 3x + 5\par\vspace{2pt}
\textbf{\color{primary} Câu 23.} Biến đổi và tính toán theo kiến thức Chương IV: Diện tích hình phẳng giới hạn bởi parabol y = x$^2$ - 4 và trục hoành Ox (trên đoạn [-2; 2]) bằng. Kết quả thu được là:\par
\hspace*{1.5em}\textbf{A.} -32/3\quad\quad \textbf{B.} 16/3\quad\quad \textbf{C.} 8\quad\quad \textbf{D.} 32/3\par\vspace{2pt}
\textbf{\color{primary} Câu 24.} Xét bài toán mức độ thông hiểu sau: Cho $\int_{1}^{3}$ f(x) dx = 7 và $\int_{1}^{3}$ g(x) dx = 2. Giá trị của tích phân I = $\int_{1}^{3}$ [2f(x) - 3g(x)] dx bằng. Đáp án đúng là:\par
\hspace*{1.5em}\textbf{A.} 5\quad\quad \textbf{B.} 12\quad\quad \textbf{C.} 20\quad\quad \textbf{D.} 8\par\vspace{2pt}
\textbf{\color{primary} Câu 25.} Để hoàn thành phiếu học tập Chương IV, hãy xác định: Tìm nguyên hàm F(x) của hàm số f(x) = 2x + 6 biết F(0) = 5:\par
\hspace*{1.5em}\textbf{A.} F(x) = x$^2$ + 5\quad\quad \textbf{B.} F(x) = x$^2$ + 6x + 5\quad\quad \textbf{C.} F(x) = 2x$^2$ + 6x + 5\quad\quad \textbf{D.} F(x) = x$^2$ + 6x - 5\par\vspace{2pt}
\textbf{\color{primary} Câu 26.} Phân tích dữ kiện bài toán: Diện tích hình phẳng giới hạn bởi parabol y = x$^2$ - 4 và trục hoành Ox (trên đoạn [-2; 2]) bằng. Giá trị cần tìm bằng:\par
\hspace*{1.5em}\textbf{A.} 8\quad\quad \textbf{B.} 16/3\quad\quad \textbf{C.} -32/3\quad\quad \textbf{D.} 32/3\par\vspace{2pt}
\textbf{\color{primary} Câu 27.} Áp dụng tính chất đã học ở Chương IV vào bài tập sau: Cho $\int_{1}^{3}$ f(x) dx = 10 và $\int_{1}^{3}$ g(x) dx = 2. Giá trị của tích phân I = $\int_{1}^{3}$ [2f(x) - 3g(x)] dx bằng:\par
\hspace*{1.5em}\textbf{A.} 14\quad\quad \textbf{B.} 18\quad\quad \textbf{C.} 8\quad\quad \textbf{D.} 26\par\vspace{2pt}
\textbf{\color{primary} Câu 28.} Thực hiện các bước suy luận hợp lý để giải quyết yêu cầu: Tìm nguyên hàm F(x) của hàm số f(x) = 2x + 9 biết F(0) = 5:\par
\hspace*{1.5em}\textbf{A.} F(x) = x$^2$ + 5\quad\quad \textbf{B.} F(x) = 2x$^2$ + 9x + 5\quad\quad \textbf{C.} F(x) = x$^2$ + 9x + 5\quad\quad \textbf{D.} F(x) = x$^2$ + 9x - 5\par\vspace{2pt}
\textbf{\color{primary} Câu 29.} Cho bài toán tự luyện (Nguyên hàm và tích phân): Diện tích hình phẳng giới hạn bởi parabol y = x$^2$ - 4 và trục hoành Ox (trên đoạn [-2; 2]) bằng. Phương án nào sau đây chính xác?\par
\hspace*{1.5em}\textbf{A.} 16/3\quad\quad \textbf{B.} -32/3\quad\quad \textbf{C.} 32/3\quad\quad \textbf{D.} 8\par\vspace{2pt}
\textbf{\color{primary} Câu 30.} Kết hợp các phép biến đổi trong Chương IV, hãy tìm kết quả của bài toán: Cho $\int_{1}^{3}$ f(x) dx = 13 và $\int_{1}^{3}$ g(x) dx = 2. Giá trị của tích phân I = $\int_{1}^{3}$ [2f(x) - 3g(x)] dx bằng:\par
\hspace*{1.5em}\textbf{A.} 11\quad\quad \textbf{B.} 24\quad\quad \textbf{C.} 20\quad\quad \textbf{D.} 32\par\vspace{2pt}
\textbf{\color{primary} Câu 31.} Một bài kiểm tra định kỳ nội dung Chương IV nêu câu hỏi: Tìm nguyên hàm F(x) của hàm số f(x) = 2x + 12 biết F(0) = 5. Hãy chọn đáp số đúng:\par
\hspace*{1.5em}\textbf{A.} F(x) = 2x$^2$ + 12x + 5\quad\quad \textbf{B.} F(x) = x$^2$ + 5\quad\quad \textbf{C.} F(x) = x$^2$ + 12x - 5\quad\quad \textbf{D.} F(x) = x$^2$ + 12x + 5\par\vspace{2pt}
\textbf{\color{primary} Câu 32.} Tìm lời giải đúng cho bài toán thông hiểu: Diện tích hình phẳng giới hạn bởi parabol y = x$^2$ - 4 và trục hoành Ox (trên đoạn [-2; 2]) bằng:\par
\hspace*{1.5em}\textbf{A.} 8\quad\quad \textbf{B.} 32/3\quad\quad \textbf{C.} -32/3\quad\quad \textbf{D.} 16/3\par\vspace{2pt}
\textbf{\color{primary} Câu 33.} Sử dụng mối liên hệ giữa các đại lượng trong Chương IV: Cho $\int_{1}^{3}$ f(x) dx = 16 và $\int_{1}^{3}$ g(x) dx = 2. Giá trị của tích phân I = $\int_{1}^{3}$ [2f(x) - 3g(x)] dx bằng:\par
\hspace*{1.5em}\textbf{A.} 14\quad\quad \textbf{B.} 38\quad\quad \textbf{C.} 30\quad\quad \textbf{D.} 26\par\vspace{2pt}
\textbf{\color{primary} Câu 34.} Hãy tính toán và chọn phương án đúng đối với bài tập: Tìm nguyên hàm F(x) của hàm số f(x) = 2x + 15 biết F(0) = 5:\par
\hspace*{1.5em}\textbf{A.} F(x) = x$^2$ + 15x - 5\quad\quad \textbf{B.} F(x) = 2x$^2$ + 15x + 5\quad\quad \textbf{C.} F(x) = x$^2$ + 15x + 5\quad\quad \textbf{D.} F(x) = x$^2$ + 5\par\vspace{2pt}
\textbf{\color{primary} Câu 35.} Dựa trên các tính chất trọng tâm của Chương IV, xét bài toán: Diện tích hình phẳng giới hạn bởi parabol y = x$^2$ - 4 và trục hoành Ox (trên đoạn [-2; 2]) bằng:\par
\hspace*{1.5em}\textbf{A.} 16/3\quad\quad \textbf{B.} 32/3\quad\quad \textbf{C.} -32/3\quad\quad \textbf{D.} 8\par\vspace{2pt}
\textbf{\color{primary} Câu 36.} Thực hành giải toán (Nguyên hàm và tích phân): Cho $\int_{1}^{3}$ f(x) dx = 19 và $\int_{1}^{3}$ g(x) dx = 2. Giá trị của tích phân I = $\int_{1}^{3}$ [2f(x) - 3g(x)] dx bằng. Đáp án nào dưới đây thỏa mãn?\par
\hspace*{1.5em}\textbf{A.} 36\quad\quad \textbf{B.} 44\quad\quad \textbf{C.} 17\quad\quad \textbf{D.} 32\par\vspace{2pt}
\textbf{\color{primary} Câu 37.} Kiểm tra kỹ năng biến đổi và suy luận qua bài toán: Tìm nguyên hàm F(x) của hàm số f(x) = 2x + 18 biết F(0) = 5:\par
\hspace*{1.5em}\textbf{A.} F(x) = 2x$^2$ + 18x + 5\quad\quad \textbf{B.} F(x) = x$^2$ + 18x - 5\quad\quad \textbf{C.} F(x) = x$^2$ + 5\quad\quad \textbf{D.} F(x) = x$^2$ + 18x + 5\par\vspace{2pt}
\textbf{\color{primary} Câu 38.} Cho giả thiết thuộc phạm vi Chương IV: Diện tích hình phẳng giới hạn bởi parabol y = x$^2$ - 4 và trục hoành Ox (trên đoạn [-2; 2]) bằng. Kết luận đúng là:\par
\hspace*{1.5em}\textbf{A.} -32/3\quad\quad \textbf{B.} 8\quad\quad \textbf{C.} 16/3\quad\quad \textbf{D.} 32/3\par\vspace{2pt}
\textbf{\color{primary} Câu 39.} Rút gọn hoặc tính giá trị theo yêu cầu sau của Chương IV: Cho $\int_{1}^{3}$ f(x) dx = 22 và $\int_{1}^{3}$ g(x) dx = 2. Giá trị của tích phân I = $\int_{1}^{3}$ [2f(x) - 3g(x)] dx bằng:\par
\hspace*{1.5em}\textbf{A.} 38\quad\quad \textbf{B.} 20\quad\quad \textbf{C.} 50\quad\quad \textbf{D.} 42\par\vspace{2pt}
\textbf{\color{primary} Câu 40.} Xác định phương án đúng sau khi thực hiện phép tính / suy luận: Tìm nguyên hàm F(x) của hàm số f(x) = 2x + 21 biết F(0) = 5:\par
\hspace*{1.5em}\textbf{A.} F(x) = x$^2$ + 5\quad\quad \textbf{B.} F(x) = 2x$^2$ + 21x + 5\quad\quad \textbf{C.} F(x) = x$^2$ + 21x - 5\quad\quad \textbf{D.} F(x) = x$^2$ + 21x + 5\par\vspace{2pt}
\subsection*{\color{accent} Phần III. Mức độ Khó (Vận dụng \& Vận dụng cao -- Câu 41 đến Câu 50)}
\textbf{\color{primary} Câu 41.} Thể tích khối tròn xoay tạo thành khi quay hình phẳng giới hạn bởi đồ thị hàm số y = $\sqrt{x}$, trục hoành Ox và hai đường thẳng x = 0, x = 4 quanh trục Ox bằng:\par
\hspace*{1.5em}\textbf{A.} 4$\pi$\quad\quad \textbf{B.} 8$\pi$\quad\quad \textbf{C.} 8\quad\quad \textbf{D.} 16$\pi$/3\par\vspace{2pt}
\textbf{\color{primary} Câu 42.} Bài toán vận dụng (Chương IV – Nguyên hàm và tích phân): Một vật chuyển động với vận tốc v(t) = 2t + 3 (m/s). Quãng đường vật đi được từ thời điểm t = 0 đến thời điểm t = 3 (giây) bằng. Kết quả đúng là:\par
\hspace*{1.5em}\textbf{A.} 9 m\quad\quad \textbf{B.} 18 m\quad\quad \textbf{C.} 27 m\quad\quad \textbf{D.} 12 m\par\vspace{2pt}
\textbf{\color{primary} Câu 43.} Vận dụng kiến thức nâng cao trong phạm vi Chương IV để giải bài toán: Thể tích khối tròn xoay tạo thành khi quay hình phẳng giới hạn bởi đồ thị hàm số y = $\sqrt{x}$, trục hoành Ox và hai đường thẳng x = 0, x = 4 quanh trục Ox bằng:\par
\hspace*{1.5em}\textbf{A.} 8\quad\quad \textbf{B.} 16$\pi$/3\quad\quad \textbf{C.} 4$\pi$\quad\quad \textbf{D.} 8$\pi$\par\vspace{2pt}
\textbf{\color{primary} Câu 44.} Phân tích và giải bài toán tổng hợp sau của Chương IV: Một vật chuyển động với vận tốc v(t) = 2t + 5 (m/s). Quãng đường vật đi được từ thời điểm t = 0 đến thời điểm t = 3 (giây) bằng:\par
\hspace*{1.5em}\textbf{A.} 11 m\quad\quad \textbf{B.} 33 m\quad\quad \textbf{C.} 24 m\quad\quad \textbf{D.} 14 m\par\vspace{2pt}
\textbf{\color{primary} Câu 45.} Thử thách tư duy trong phạm vi kiến thức Chương IV: Thể tích khối tròn xoay tạo thành khi quay hình phẳng giới hạn bởi đồ thị hàm số y = $\sqrt{x}$, trục hoành Ox và hai đường thẳng x = 0, x = 4 quanh trục Ox bằng. Đáp số chính xác là:\par
\hspace*{1.5em}\textbf{A.} 8\quad\quad \textbf{B.} 4$\pi$\quad\quad \textbf{C.} 8$\pi$\quad\quad \textbf{D.} 16$\pi$/3\par\vspace{2pt}
\textbf{\color{primary} Câu 46.} Sử dụng phương pháp giải toán của Chương IV để tìm đáp án cho bài toán: Một vật chuyển động với vận tốc v(t) = 2t + 7 (m/s). Quãng đường vật đi được từ thời điểm t = 0 đến thời điểm t = 3 (giây) bằng:\par
\hspace*{1.5em}\textbf{A.} 30 m\quad\quad \textbf{B.} 13 m\quad\quad \textbf{C.} 16 m\quad\quad \textbf{D.} 39 m\par\vspace{2pt}
\textbf{\color{primary} Câu 47.} Xét bài toán vận dụng thực tiễn / nâng cao thuộc Chương IV: Thể tích khối tròn xoay tạo thành khi quay hình phẳng giới hạn bởi đồ thị hàm số y = $\sqrt{x}$, trục hoành Ox và hai đường thẳng x = 0, x = 4 quanh trục Ox bằng:\par
\hspace*{1.5em}\textbf{A.} 16$\pi$/3\quad\quad \textbf{B.} 8$\pi$\quad\quad \textbf{C.} 8\quad\quad \textbf{D.} 4$\pi$\par\vspace{2pt}
\textbf{\color{primary} Câu 48.} Tìm kết quả cuối cùng của bài toán phân loại học sinh khá giỏi (Chương IV): Một vật chuyển động với vận tốc v(t) = 2t + 9 (m/s). Quãng đường vật đi được từ thời điểm t = 0 đến thời điểm t = 3 (giây) bằng:\par
\hspace*{1.5em}\textbf{A.} 36 m\quad\quad \textbf{B.} 45 m\quad\quad \textbf{C.} 18 m\quad\quad \textbf{D.} 15 m\par\vspace{2pt}
\textbf{\color{primary} Câu 49.} Kết hợp các tính chất đã học trong Chương IV, hãy giải quyết bài toán: Thể tích khối tròn xoay tạo thành khi quay hình phẳng giới hạn bởi đồ thị hàm số y = $\sqrt{x}$, trục hoành Ox và hai đường thẳng x = 0, x = 4 quanh trục Ox bằng:\par
\hspace*{1.5em}\textbf{A.} 8\quad\quad \textbf{B.} 8$\pi$\quad\quad \textbf{C.} 16$\pi$/3\quad\quad \textbf{D.} 4$\pi$\par\vspace{2pt}
\textbf{\color{primary} Câu 50.} Câu hỏi vận dụng cao cuối Chương IV: Một vật chuyển động với vận tốc v(t) = 2t + 11 (m/s). Quãng đường vật đi được từ thời điểm t = 0 đến thời điểm t = 3 (giây) bằng. Phương án đúng là:\par
\hspace*{1.5em}\textbf{A.} 51 m\quad\quad \textbf{B.} 17 m\quad\quad \textbf{C.} 42 m\quad\quad \textbf{D.} 20 m\par\vspace{2pt}
\subsection*{\color{primary} Bảng đáp án Chương IV: Nguyên hàm và tích phân}
\begin{center}
\renewcommand{\arraystretch}{1.25}
\begin{tabularx}{\textwidth}{*{10}{>{\centering\arraybackslash}X}}
\toprule
\textbf{\color{primary} Cột 1} & \textbf{\color{primary} Cột 2} & \textbf{\color{primary} Cột 3} & \textbf{\color{primary} Cột 4} & \textbf{\color{primary} Cột 5} & \textbf{\color{primary} Cột 6} & \textbf{\color{primary} Cột 7} & \textbf{\color{primary} Cột 8} & \textbf{\color{primary} Cột 9} & \textbf{\color{primary} Cột 10} \\
\midrule
\textbf{1.C} & \textbf{2.B} & \textbf{3.A} & \textbf{4.A} & \textbf{5.A} & \textbf{6.C} & \textbf{7.B} & \textbf{8.D} & \textbf{9.C} & \textbf{10.D} \\
\rowcolor{rowtint} \textbf{11.C} & \textbf{12.A} & \textbf{13.B} & \textbf{14.D} & \textbf{15.A} & \textbf{16.B} & \textbf{17.A} & \textbf{18.B} & \textbf{19.B} & \textbf{20.B} \\
\textbf{21.A} & \textbf{22.D} & \textbf{23.D} & \textbf{24.D} & \textbf{25.B} & \textbf{26.D} & \textbf{27.A} & \textbf{28.C} & \textbf{29.C} & \textbf{30.C} \\
\rowcolor{rowtint} \textbf{31.D} & \textbf{32.B} & \textbf{33.D} & \textbf{34.C} & \textbf{35.B} & \textbf{36.D} & \textbf{37.D} & \textbf{38.D} & \textbf{39.A} & \textbf{40.D} \\
\textbf{41.B} & \textbf{42.B} & \textbf{43.D} & \textbf{44.C} & \textbf{45.C} & \textbf{46.A} & \textbf{47.B} & \textbf{48.A} & \textbf{49.B} & \textbf{50.C} \\
\bottomrule
\end{tabularx}
\end{center}\vspace{10pt}
\section*{\color{primary} CHƯƠNG V: PHƯƠNG PHÁP TỌA ĐỘ TRONG KHÔNG GIAN}
\textit{Trọng tâm kiến thức SGK: Phương trình mặt phẳng, phương trình đường thẳng trong không gian (tham số, chính tắc), công thức tính góc, khoảng cách và phương trình mặt cầu.}\par\vspace{4pt}
\subsection*{\color{accent} Phần I. Mức độ Dễ (Nhận biết -- Câu 1 đến Câu 20)}
\textbf{\color{primary} Câu 1.} Trong không gian Oxyz, mặt phẳng (P): 2x - 3y + 4z - 5 = 0 có một vectơ pháp tuyến là:\par
\hspace*{1.5em}\textbf{A.} $\vec{n}$ = (2; 4; -5)\quad\quad \textbf{B.} $\vec{n}$ = (2; 3; 4)\quad\quad \textbf{C.} $\vec{n}$ = (-2; -3; 4)\quad\quad \textbf{D.} $\vec{n}$ = (2; -3; 4)\par\vspace{2pt}
\textbf{\color{primary} Câu 2.} Thực hiện yêu cầu nhận biết sau: Trong không gian Oxyz, đường thẳng d: (x - 1)/2 = (y + 2)/(-3) = (z - 2)/5 có một vectơ chỉ phương là.\par
\hspace*{1.5em}\textbf{A.} $\vec{u}$ = (2; -3; 5)\quad\quad \textbf{B.} $\vec{u}$ = (2; 3; 5)\quad\quad \textbf{C.} $\vec{u}$ = (1; -2; 2)\quad\quad \textbf{D.} $\vec{u}$ = (-1; 2; -2)\par\vspace{2pt}
\textbf{\color{primary} Câu 3.} Trong bài học thuộc Chương V, xét nội dung: Trong không gian Oxyz, tâm I và bán kính R của mặt cầu (S): (x - 1)$^2$ + (y + 2)$^2$ + (z - 3)$^2$ = 25 là. Đáp án đúng là:\par
\hspace*{1.5em}\textbf{A.} I(1; -2; 3), R = 25\quad\quad \textbf{B.} I(-1; 2; -3), R = 5\quad\quad \textbf{C.} I(1; 2; 3), R = 5\quad\quad \textbf{D.} I(1; -2; 3), R = 5\par\vspace{2pt}
\textbf{\color{primary} Câu 4.} Dựa vào kiến thức chuẩn SGK của Chương V (Phương pháp tọa độ trong không gian), trong không gian Oxyz, mặt phẳng (P): 2x - 6y + 4z - 5 = 0 có một vectơ pháp tuyến là:\par
\hspace*{1.5em}\textbf{A.} $\vec{n}$ = (2; -6; 4)\quad\quad \textbf{B.} $\vec{n}$ = (-2; -6; 4)\quad\quad \textbf{C.} $\vec{n}$ = (2; 4; -5)\quad\quad \textbf{D.} $\vec{n}$ = (2; 6; 4)\par\vspace{2pt}
\textbf{\color{primary} Câu 5.} Chọn kết quả đúng nhất khi xét: Trong không gian Oxyz, đường thẳng d: (x - 1)/2 = (y + 2)/(-3) = (z - 5)/5 có một vectơ chỉ phương là:\par
\hspace*{1.5em}\textbf{A.} $\vec{u}$ = (-1; 2; -5)\quad\quad \textbf{B.} $\vec{u}$ = (2; 3; 5)\quad\quad \textbf{C.} $\vec{u}$ = (2; -3; 5)\quad\quad \textbf{D.} $\vec{u}$ = (1; -2; 5)\par\vspace{2pt}
\textbf{\color{primary} Câu 6.} Khi làm bài tập cơ bản về phương pháp tọa độ trong không gian, một học sinh gặp câu hỏi: Trong không gian Oxyz, tâm I và bán kính R của mặt cầu (S): (x - 1)$^2$ + (y + 2)$^2$ + (z - 3)$^2$ = 64 là. Kết quả nào dưới đây đúng?\par
\hspace*{1.5em}\textbf{A.} I(1; 2; 3), R = 8\quad\quad \textbf{B.} I(1; -2; 3), R = 64\quad\quad \textbf{C.} I(1; -2; 3), R = 8\quad\quad \textbf{D.} I(-1; 2; -3), R = 8\par\vspace{2pt}
\textbf{\color{primary} Câu 7.} Xác định giá trị hoặc khẳng định đúng cho bài toán sau: Trong không gian Oxyz, mặt phẳng (P): 2x - 9y + 4z - 5 = 0 có một vectơ pháp tuyến là:\par
\hspace*{1.5em}\textbf{A.} $\vec{n}$ = (2; 9; 4)\quad\quad \textbf{B.} $\vec{n}$ = (2; -9; 4)\quad\quad \textbf{C.} $\vec{n}$ = (2; 4; -5)\quad\quad \textbf{D.} $\vec{n}$ = (-2; -9; 4)\par\vspace{2pt}
\textbf{\color{primary} Câu 8.} Áp dụng trực tiếp định nghĩa và quy tắc trong Chương V: Trong không gian Oxyz, đường thẳng d: (x - 1)/2 = (y + 2)/(-3) = (z - 8)/5 có một vectơ chỉ phương là:\par
\hspace*{1.5em}\textbf{A.} $\vec{u}$ = (1; -2; 8)\quad\quad \textbf{B.} $\vec{u}$ = (2; -3; 5)\quad\quad \textbf{C.} $\vec{u}$ = (-1; 2; -8)\quad\quad \textbf{D.} $\vec{u}$ = (2; 3; 5)\par\vspace{2pt}
\textbf{\color{primary} Câu 9.} Kiểm tra nhanh kiến thức nền tảng: Trong không gian Oxyz, tâm I và bán kính R của mặt cầu (S): (x - 1)$^2$ + (y + 2)$^2$ + (z - 3)$^2$ = 121 là:\par
\hspace*{1.5em}\textbf{A.} I(1; -2; 3), R = 121\quad\quad \textbf{B.} I(1; 2; 3), R = 11\quad\quad \textbf{C.} I(1; -2; 3), R = 11\quad\quad \textbf{D.} I(-1; 2; -3), R = 11\par\vspace{2pt}
\textbf{\color{primary} Câu 10.} Cho dữ kiện nhận biết sau đây: Trong không gian Oxyz, mặt phẳng (P): 2x - 12y + 4z - 5 = 0 có một vectơ pháp tuyến là. Phương án chính xác là:\par
\hspace*{1.5em}\textbf{A.} $\vec{n}$ = (2; -12; 4)\quad\quad \textbf{B.} $\vec{n}$ = (2; 4; -5)\quad\quad \textbf{C.} $\vec{n}$ = (-2; -12; 4)\quad\quad \textbf{D.} $\vec{n}$ = (2; 12; 4)\par\vspace{2pt}
\textbf{\color{primary} Câu 11.} Hãy chỉ ra kết quả đúng trong bốn lựa chọn A, B, C, D đối với yêu cầu: Trong không gian Oxyz, đường thẳng d: (x - 1)/2 = (y + 2)/(-3) = (z - 11)/5 có một vectơ chỉ phương là:\par
\hspace*{1.5em}\textbf{A.} $\vec{u}$ = (-1; 2; -11)\quad\quad \textbf{B.} $\vec{u}$ = (2; -3; 5)\quad\quad \textbf{C.} $\vec{u}$ = (1; -2; 11)\quad\quad \textbf{D.} $\vec{u}$ = (2; 3; 5)\par\vspace{2pt}
\textbf{\color{primary} Câu 12.} Câu hỏi ôn tập lý thuyết và kỹ năng cơ bản (Chương V): Trong không gian Oxyz, tâm I và bán kính R của mặt cầu (S): (x - 1)$^2$ + (y + 2)$^2$ + (z - 3)$^2$ = 196 là:\par
\hspace*{1.5em}\textbf{A.} I(1; 2; 3), R = 14\quad\quad \textbf{B.} I(-1; 2; -3), R = 14\quad\quad \textbf{C.} I(1; -2; 3), R = 196\quad\quad \textbf{D.} I(1; -2; 3), R = 14\par\vspace{2pt}
\textbf{\color{primary} Câu 13.} Theo nội dung sách giáo khoa Kết nối tri thức với cuộc sống (Chương V), trong không gian Oxyz, mặt phẳng (P): 2x - 15y + 4z - 5 = 0 có một vectơ pháp tuyến là:\par
\hspace*{1.5em}\textbf{A.} $\vec{n}$ = (2; 15; 4)\quad\quad \textbf{B.} $\vec{n}$ = (2; 4; -5)\quad\quad \textbf{C.} $\vec{n}$ = (-2; -15; 4)\quad\quad \textbf{D.} $\vec{n}$ = (2; -15; 4)\par\vspace{2pt}
\textbf{\color{primary} Câu 14.} Tìm đáp số đúng của bài toán cơ bản sau: Trong không gian Oxyz, đường thẳng d: (x - 1)/2 = (y + 2)/(-3) = (z - 14)/5 có một vectơ chỉ phương là:\par
\hspace*{1.5em}\textbf{A.} $\vec{u}$ = (-1; 2; -14)\quad\quad \textbf{B.} $\vec{u}$ = (2; 3; 5)\quad\quad \textbf{C.} $\vec{u}$ = (2; -3; 5)\quad\quad \textbf{D.} $\vec{u}$ = (1; -2; 14)\par\vspace{2pt}
\textbf{\color{primary} Câu 15.} Xét đối tượng toán học trong Chương V: Trong không gian Oxyz, tâm I và bán kính R của mặt cầu (S): (x - 1)$^2$ + (y + 2)$^2$ + (z - 3)$^2$ = 289 là. Giá trị hoặc tính chất tương ứng là:\par
\hspace*{1.5em}\textbf{A.} I(-1; 2; -3), R = 17\quad\quad \textbf{B.} I(1; -2; 3), R = 289\quad\quad \textbf{C.} I(1; -2; 3), R = 17\quad\quad \textbf{D.} I(1; 2; 3), R = 17\par\vspace{2pt}
\textbf{\color{primary} Câu 16.} Hoàn thành phép xác định cơ bản sau đây: Trong không gian Oxyz, mặt phẳng (P): 2x - 18y + 4z - 5 = 0 có một vectơ pháp tuyến là:\par
\hspace*{1.5em}\textbf{A.} $\vec{n}$ = (2; 4; -5)\quad\quad \textbf{B.} $\vec{n}$ = (2; -18; 4)\quad\quad \textbf{C.} $\vec{n}$ = (2; 18; 4)\quad\quad \textbf{D.} $\vec{n}$ = (-2; -18; 4)\par\vspace{2pt}
\textbf{\color{primary} Câu 17.} Trong các khẳng định hoặc giá trị dưới đây, đâu là đáp án đúng cho câu hỏi: Trong không gian Oxyz, đường thẳng d: (x - 1)/2 = (y + 2)/(-3) = (z - 17)/5 có một vectơ chỉ phương là?\par
\hspace*{1.5em}\textbf{A.} $\vec{u}$ = (1; -2; 17)\quad\quad \textbf{B.} $\vec{u}$ = (-1; 2; -17)\quad\quad \textbf{C.} $\vec{u}$ = (2; -3; 5)\quad\quad \textbf{D.} $\vec{u}$ = (2; 3; 5)\par\vspace{2pt}
\textbf{\color{primary} Câu 18.} Bài tập khởi động Chương V: Trong không gian Oxyz, tâm I và bán kính R của mặt cầu (S): (x - 1)$^2$ + (y + 2)$^2$ + (z - 3)$^2$ = 400 là:\par
\hspace*{1.5em}\textbf{A.} I(-1; 2; -3), R = 20\quad\quad \textbf{B.} I(1; -2; 3), R = 20\quad\quad \textbf{C.} I(1; 2; 3), R = 20\quad\quad \textbf{D.} I(1; -2; 3), R = 400\par\vspace{2pt}
\textbf{\color{primary} Câu 19.} Quan sát dữ kiện đề bài: Trong không gian Oxyz, mặt phẳng (P): 2x - 21y + 4z - 5 = 0 có một vectơ pháp tuyến là. Hãy chọn phương án đúng:\par
\hspace*{1.5em}\textbf{A.} $\vec{n}$ = (2; -21; 4)\quad\quad \textbf{B.} $\vec{n}$ = (2; 4; -5)\quad\quad \textbf{C.} $\vec{n}$ = (-2; -21; 4)\quad\quad \textbf{D.} $\vec{n}$ = (2; 21; 4)\par\vspace{2pt}
\textbf{\color{primary} Câu 20.} Sử dụng công thức hoặc tính chất đầu bài của Chương V để giải quyết yêu cầu: Trong không gian Oxyz, đường thẳng d: (x - 1)/2 = (y + 2)/(-3) = (z - 20)/5 có một vectơ chỉ phương là:\par
\hspace*{1.5em}\textbf{A.} $\vec{u}$ = (2; -3; 5)\quad\quad \textbf{B.} $\vec{u}$ = (-1; 2; -20)\quad\quad \textbf{C.} $\vec{u}$ = (2; 3; 5)\quad\quad \textbf{D.} $\vec{u}$ = (1; -2; 20)\par\vspace{2pt}
\subsection*{\color{accent} Phần II. Mức độ Trung bình (Thông hiểu -- Câu 21 đến Câu 40)}
\textbf{\color{primary} Câu 21.} Trong không gian Oxyz, khoảng cách từ điểm M(1; 2; 2) đến mặt phẳng (P): 2x + 2y + z + -5 = 0 bằng:\par
\hspace*{1.5em}\textbf{A.} 2\quad\quad \textbf{B.} 3\quad\quad \textbf{C.} 1\quad\quad \textbf{D.} 0.33\par\vspace{2pt}
\textbf{\color{primary} Câu 22.} Vận dụng quy tắc trong Chương V (Phương pháp tọa độ trong không gian) để giải bài toán: Trong không gian Oxyz, phương trình mặt phẳng đi qua điểm A(1; 0; 2) và có vectơ pháp tuyến $\vec{n}$ = (2; -1; 4) là:\par
\hspace*{1.5em}\textbf{A.} 2x + y + 4z - 5 = 0\quad\quad \textbf{B.} x + 2z - 10 = 0\quad\quad \textbf{C.} 2x - y + 4z - 10 = 0\quad\quad \textbf{D.} 2x - y + 4z + 10 = 0\par\vspace{2pt}
\textbf{\color{primary} Câu 23.} Biến đổi và tính toán theo kiến thức Chương V: Trong không gian Oxyz, phương trình mặt cầu tâm I(1; -2; 0) và đi qua điểm A(3; 0; 1) là. Kết quả thu được là:\par
\begin{itemize}\setlength{\itemsep}{1pt}\setlength{\parskip}{0pt}
  \item[] \textbf{A.} (x + 1)$^2$ + (y - 2)$^2$ + z$^2$ = 9
  \item[] \textbf{B.} (x - 1)$^2$ + (y + 2)$^2$ + z$^2$ = 9
  \item[] \textbf{C.} (x - 1)$^2$ + (y + 2)$^2$ + z$^2$ = 3
  \item[] \textbf{D.} (x - 3)$^2$ + y$^2$ + (z - 1)$^2$ = 9
\end{itemize}\vspace{2pt}
\textbf{\color{primary} Câu 24.} Xét bài toán mức độ thông hiểu sau: Trong không gian Oxyz, khoảng cách từ điểm M(1; 2; 2) đến mặt phẳng (P): 2x + 2y + z + 4 = 0 bằng. Đáp án đúng là:\par
\hspace*{1.5em}\textbf{A.} 4\quad\quad \textbf{B.} 5\quad\quad \textbf{C.} 1.33\quad\quad \textbf{D.} 12\par\vspace{2pt}
\textbf{\color{primary} Câu 25.} Để hoàn thành phiếu học tập Chương V, hãy xác định: Trong không gian Oxyz, phương trình mặt phẳng đi qua điểm A(1; 0; 2) và có vectơ pháp tuyến $\vec{n}$ = (2; -1; 7) là:\par
\hspace*{1.5em}\textbf{A.} 2x + y + 7z - 5 = 0\quad\quad \textbf{B.} x + 2z - 16 = 0\quad\quad \textbf{C.} 2x - y + 7z - 16 = 0\quad\quad \textbf{D.} 2x - y + 7z + 16 = 0\par\vspace{2pt}
\textbf{\color{primary} Câu 26.} Phân tích dữ kiện bài toán: Trong không gian Oxyz, phương trình mặt cầu tâm I(1; -2; 0) và đi qua điểm A(3; 0; 1) là. Giá trị cần tìm bằng:\par
\begin{itemize}\setlength{\itemsep}{1pt}\setlength{\parskip}{0pt}
  \item[] \textbf{A.} (x + 1)$^2$ + (y - 2)$^2$ + z$^2$ = 9
  \item[] \textbf{B.} (x - 1)$^2$ + (y + 2)$^2$ + z$^2$ = 3
  \item[] \textbf{C.} (x - 1)$^2$ + (y + 2)$^2$ + z$^2$ = 9
  \item[] \textbf{D.} (x - 3)$^2$ + y$^2$ + (z - 1)$^2$ = 9
\end{itemize}\vspace{2pt}
\textbf{\color{primary} Câu 27.} Áp dụng tính chất đã học ở Chương V vào bài tập sau: Trong không gian Oxyz, khoảng cách từ điểm M(1; 2; 2) đến mặt phẳng (P): 2x + 2y + z + 13 = 0 bằng:\par
\hspace*{1.5em}\textbf{A.} 7\quad\quad \textbf{B.} 21\quad\quad \textbf{C.} 8\quad\quad \textbf{D.} 2.33\par\vspace{2pt}
\textbf{\color{primary} Câu 28.} Thực hiện các bước suy luận hợp lý để giải quyết yêu cầu: Trong không gian Oxyz, phương trình mặt phẳng đi qua điểm A(1; 0; 2) và có vectơ pháp tuyến $\vec{n}$ = (2; -1; 10) là:\par
\hspace*{1.5em}\textbf{A.} x + 2z - 22 = 0\quad\quad \textbf{B.} 2x - y + 10z - 22 = 0\quad\quad \textbf{C.} 2x + y + 10z - 5 = 0\quad\quad \textbf{D.} 2x - y + 10z + 22 = 0\par\vspace{2pt}
\textbf{\color{primary} Câu 29.} Cho bài toán tự luyện (Phương pháp tọa độ trong không gian): Trong không gian Oxyz, phương trình mặt cầu tâm I(1; -2; 0) và đi qua điểm A(3; 0; 1) là. Phương án nào sau đây chính xác?\par
\begin{itemize}\setlength{\itemsep}{1pt}\setlength{\parskip}{0pt}
  \item[] \textbf{A.} (x - 1)$^2$ + (y + 2)$^2$ + z$^2$ = 9
  \item[] \textbf{B.} (x + 1)$^2$ + (y - 2)$^2$ + z$^2$ = 9
  \item[] \textbf{C.} (x - 3)$^2$ + y$^2$ + (z - 1)$^2$ = 9
  \item[] \textbf{D.} (x - 1)$^2$ + (y + 2)$^2$ + z$^2$ = 3
\end{itemize}\vspace{2pt}
\textbf{\color{primary} Câu 30.} Kết hợp các phép biến đổi trong Chương V, hãy tìm kết quả của bài toán: Trong không gian Oxyz, khoảng cách từ điểm M(1; 2; 2) đến mặt phẳng (P): 2x + 2y + z + 22 = 0 bằng:\par
\hspace*{1.5em}\textbf{A.} 30\quad\quad \textbf{B.} 3.33\quad\quad \textbf{C.} 10\quad\quad \textbf{D.} 11\par\vspace{2pt}
\textbf{\color{primary} Câu 31.} Một bài kiểm tra định kỳ nội dung Chương V nêu câu hỏi: Trong không gian Oxyz, phương trình mặt phẳng đi qua điểm A(1; 0; 2) và có vectơ pháp tuyến $\vec{n}$ = (2; -1; 13) là. Hãy chọn đáp số đúng:\par
\hspace*{1.5em}\textbf{A.} x + 2z - 28 = 0\quad\quad \textbf{B.} 2x + y + 13z - 5 = 0\quad\quad \textbf{C.} 2x - y + 13z + 28 = 0\quad\quad \textbf{D.} 2x - y + 13z - 28 = 0\par\vspace{2pt}
\textbf{\color{primary} Câu 32.} Tìm lời giải đúng cho bài toán thông hiểu: Trong không gian Oxyz, phương trình mặt cầu tâm I(1; -2; 0) và đi qua điểm A(3; 0; 1) là:\par
\begin{itemize}\setlength{\itemsep}{1pt}\setlength{\parskip}{0pt}
  \item[] \textbf{A.} (x + 1)$^2$ + (y - 2)$^2$ + z$^2$ = 9
  \item[] \textbf{B.} (x - 1)$^2$ + (y + 2)$^2$ + z$^2$ = 3
  \item[] \textbf{C.} (x - 1)$^2$ + (y + 2)$^2$ + z$^2$ = 9
  \item[] \textbf{D.} (x - 3)$^2$ + y$^2$ + (z - 1)$^2$ = 9
\end{itemize}\vspace{2pt}
\textbf{\color{primary} Câu 33.} Sử dụng mối liên hệ giữa các đại lượng trong Chương V: Trong không gian Oxyz, khoảng cách từ điểm M(1; 2; 2) đến mặt phẳng (P): 2x + 2y + z + 31 = 0 bằng:\par
\hspace*{1.5em}\textbf{A.} 14\quad\quad \textbf{B.} 4.33\quad\quad \textbf{C.} 13\quad\quad \textbf{D.} 39\par\vspace{2pt}
\textbf{\color{primary} Câu 34.} Hãy tính toán và chọn phương án đúng đối với bài tập: Trong không gian Oxyz, phương trình mặt phẳng đi qua điểm A(1; 0; 2) và có vectơ pháp tuyến $\vec{n}$ = (2; -1; 16) là:\par
\hspace*{1.5em}\textbf{A.} 2x - y + 16z - 34 = 0\quad\quad \textbf{B.} 2x + y + 16z - 5 = 0\quad\quad \textbf{C.} x + 2z - 34 = 0\quad\quad \textbf{D.} 2x - y + 16z + 34 = 0\par\vspace{2pt}
\textbf{\color{primary} Câu 35.} Dựa trên các tính chất trọng tâm của Chương V, xét bài toán: Trong không gian Oxyz, phương trình mặt cầu tâm I(1; -2; 0) và đi qua điểm A(3; 0; 1) là:\par
\begin{itemize}\setlength{\itemsep}{1pt}\setlength{\parskip}{0pt}
  \item[] \textbf{A.} (x - 1)$^2$ + (y + 2)$^2$ + z$^2$ = 3
  \item[] \textbf{B.} (x - 3)$^2$ + y$^2$ + (z - 1)$^2$ = 9
  \item[] \textbf{C.} (x + 1)$^2$ + (y - 2)$^2$ + z$^2$ = 9
  \item[] \textbf{D.} (x - 1)$^2$ + (y + 2)$^2$ + z$^2$ = 9
\end{itemize}\vspace{2pt}
\textbf{\color{primary} Câu 36.} Thực hành giải toán (Phương pháp tọa độ trong không gian): Trong không gian Oxyz, khoảng cách từ điểm M(1; 2; 2) đến mặt phẳng (P): 2x + 2y + z + 40 = 0 bằng. Đáp án nào dưới đây thỏa mãn?\par
\hspace*{1.5em}\textbf{A.} 16\quad\quad \textbf{B.} 5.33\quad\quad \textbf{C.} 48\quad\quad \textbf{D.} 17\par\vspace{2pt}
\textbf{\color{primary} Câu 37.} Kiểm tra kỹ năng biến đổi và suy luận qua bài toán: Trong không gian Oxyz, phương trình mặt phẳng đi qua điểm A(1; 0; 2) và có vectơ pháp tuyến $\vec{n}$ = (2; -1; 19) là:\par
\hspace*{1.5em}\textbf{A.} 2x + y + 19z - 5 = 0\quad\quad \textbf{B.} 2x - y + 19z + 40 = 0\quad\quad \textbf{C.} 2x - y + 19z - 40 = 0\quad\quad \textbf{D.} x + 2z - 40 = 0\par\vspace{2pt}
\textbf{\color{primary} Câu 38.} Cho giả thiết thuộc phạm vi Chương V: Trong không gian Oxyz, phương trình mặt cầu tâm I(1; -2; 0) và đi qua điểm A(3; 0; 1) là. Kết luận đúng là:\par
\begin{itemize}\setlength{\itemsep}{1pt}\setlength{\parskip}{0pt}
  \item[] \textbf{A.} (x - 1)$^2$ + (y + 2)$^2$ + z$^2$ = 9
  \item[] \textbf{B.} (x - 1)$^2$ + (y + 2)$^2$ + z$^2$ = 3
  \item[] \textbf{C.} (x + 1)$^2$ + (y - 2)$^2$ + z$^2$ = 9
  \item[] \textbf{D.} (x - 3)$^2$ + y$^2$ + (z - 1)$^2$ = 9
\end{itemize}\vspace{2pt}
\textbf{\color{primary} Câu 39.} Rút gọn hoặc tính giá trị theo yêu cầu sau của Chương V: Trong không gian Oxyz, khoảng cách từ điểm M(1; 2; 2) đến mặt phẳng (P): 2x + 2y + z + 49 = 0 bằng:\par
\hspace*{1.5em}\textbf{A.} 20\quad\quad \textbf{B.} 19\quad\quad \textbf{C.} 6.33\quad\quad \textbf{D.} 57\par\vspace{2pt}
\textbf{\color{primary} Câu 40.} Xác định phương án đúng sau khi thực hiện phép tính / suy luận: Trong không gian Oxyz, phương trình mặt phẳng đi qua điểm A(1; 0; 2) và có vectơ pháp tuyến $\vec{n}$ = (2; -1; 22) là:\par
\hspace*{1.5em}\textbf{A.} 2x + y + 22z - 5 = 0\quad\quad \textbf{B.} 2x - y + 22z + 46 = 0\quad\quad \textbf{C.} 2x - y + 22z - 46 = 0\quad\quad \textbf{D.} x + 2z - 46 = 0\par\vspace{2pt}
\subsection*{\color{accent} Phần III. Mức độ Khó (Vận dụng \& Vận dụng cao -- Câu 41 đến Câu 50)}
\textbf{\color{primary} Câu 41.} Trong không gian Oxyz, tọa độ giao điểm M của đường thẳng d: x = 1 + t, y = 2 + 2t, z = 3 - t và mặt phẳng (P): x + y + z - 8 = 0 là:\par
\hspace*{1.5em}\textbf{A.} M(3; 6; 1)\quad\quad \textbf{B.} M(0; 0; 4)\quad\quad \textbf{C.} M(1; 2; 3)\quad\quad \textbf{D.} M(2; 4; 2)\par\vspace{2pt}
\textbf{\color{primary} Câu 42.} Bài toán vận dụng (Chương V – Phương pháp tọa độ trong không gian): Trong không gian Oxyz, sin của góc giữa đường thẳng d có vectơ chỉ phương $\vec{u}$ = (1; 2; 2) và mặt phẳng (P) có vectơ pháp tuyến $\vec{n}$ = (2; -1; 2) bằng. Kết quả đúng là:\par
\hspace*{1.5em}\textbf{A.} 4/9\quad\quad \textbf{B.} 1/3\quad\quad \textbf{C.} 2/3\quad\quad \textbf{D.} 5/9\par\vspace{2pt}
\textbf{\color{primary} Câu 43.} Vận dụng kiến thức nâng cao trong phạm vi Chương V để giải bài toán: Trong không gian Oxyz, giao điểm M của đường thẳng d: x = 1 + t, y = 2 + t, z = 3 + t và mặt phẳng (P): x + y + z - 9 = 0 có tọa độ là:\par
\hspace*{1.5em}\textbf{A.} M(3; 4; 5)\quad\quad \textbf{B.} M(0; 1; 2)\quad\quad \textbf{C.} M(2; 3; 4)\quad\quad \textbf{D.} M(1; 2; 3)\par\vspace{2pt}
\textbf{\color{primary} Câu 44.} Phân tích và giải bài toán tổng hợp sau của Chương V: Trong không gian Oxyz, tọa độ giao điểm M của đường thẳng d: x = 1 + t, y = 2 + 2t, z = 3 - t và mặt phẳng (P): x + y + z - 8 = 0 là:\par
\hspace*{1.5em}\textbf{A.} M(1; 2; 3)\quad\quad \textbf{B.} M(3; 6; 1)\quad\quad \textbf{C.} M(0; 0; 4)\quad\quad \textbf{D.} M(2; 4; 2)\par\vspace{2pt}
\textbf{\color{primary} Câu 45.} Thử thách tư duy trong phạm vi kiến thức Chương V: Trong không gian Oxyz, sin của góc giữa đường thẳng d có vectơ chỉ phương $\vec{u}$ = (1; 2; 2) và mặt phẳng (P) có vectơ pháp tuyến $\vec{n}$ = (2; -1; 2) bằng. Đáp số chính xác là:\par
\hspace*{1.5em}\textbf{A.} 5/9\quad\quad \textbf{B.} 2/3\quad\quad \textbf{C.} 1/3\quad\quad \textbf{D.} 4/9\par\vspace{2pt}
\textbf{\color{primary} Câu 46.} Sử dụng phương pháp giải toán của Chương V để tìm đáp án cho bài toán: Trong không gian Oxyz, giao điểm M của đường thẳng d: x = 1 + t, y = 2 + t, z = 6 + t và mặt phẳng (P): x + y + z - 12 = 0 có tọa độ là:\par
\hspace*{1.5em}\textbf{A.} M(1; 2; 6)\quad\quad \textbf{B.} M(3; 4; 8)\quad\quad \textbf{C.} M(2; 3; 7)\quad\quad \textbf{D.} M(0; 1; 5)\par\vspace{2pt}
\textbf{\color{primary} Câu 47.} Xét bài toán vận dụng thực tiễn / nâng cao thuộc Chương V: Trong không gian Oxyz, tọa độ giao điểm M của đường thẳng d: x = 1 + t, y = 2 + 2t, z = 3 - t và mặt phẳng (P): x + y + z - 8 = 0 là:\par
\hspace*{1.5em}\textbf{A.} M(2; 4; 2)\quad\quad \textbf{B.} M(3; 6; 1)\quad\quad \textbf{C.} M(1; 2; 3)\quad\quad \textbf{D.} M(0; 0; 4)\par\vspace{2pt}
\textbf{\color{primary} Câu 48.} Tìm kết quả cuối cùng của bài toán phân loại học sinh khá giỏi (Chương V): Trong không gian Oxyz, sin của góc giữa đường thẳng d có vectơ chỉ phương $\vec{u}$ = (1; 2; 2) và mặt phẳng (P) có vectơ pháp tuyến $\vec{n}$ = (2; -1; 2) bằng:\par
\hspace*{1.5em}\textbf{A.} 5/9\quad\quad \textbf{B.} 1/3\quad\quad \textbf{C.} 4/9\quad\quad \textbf{D.} 2/3\par\vspace{2pt}
\textbf{\color{primary} Câu 49.} Kết hợp các tính chất đã học trong Chương V, hãy giải quyết bài toán: Trong không gian Oxyz, giao điểm M của đường thẳng d: x = 1 + t, y = 2 + t, z = 9 + t và mặt phẳng (P): x + y + z - 15 = 0 có tọa độ là:\par
\hspace*{1.5em}\textbf{A.} M(2; 3; 10)\quad\quad \textbf{B.} M(3; 4; 11)\quad\quad \textbf{C.} M(1; 2; 9)\quad\quad \textbf{D.} M(0; 1; 8)\par\vspace{2pt}
\textbf{\color{primary} Câu 50.} Câu hỏi vận dụng cao cuối Chương V: Trong không gian Oxyz, tọa độ giao điểm M của đường thẳng d: x = 1 + t, y = 2 + 2t, z = 3 - t và mặt phẳng (P): x + y + z - 8 = 0 là. Phương án đúng là:\par
\hspace*{1.5em}\textbf{A.} M(1; 2; 3)\quad\quad \textbf{B.} M(3; 6; 1)\quad\quad \textbf{C.} M(2; 4; 2)\quad\quad \textbf{D.} M(0; 0; 4)\par\vspace{2pt}
\subsection*{\color{primary} Bảng đáp án Chương V: Phương pháp tọa độ trong không gian}
\begin{center}
\renewcommand{\arraystretch}{1.25}
\begin{tabularx}{\textwidth}{*{10}{>{\centering\arraybackslash}X}}
\toprule
\textbf{\color{primary} Cột 1} & \textbf{\color{primary} Cột 2} & \textbf{\color{primary} Cột 3} & \textbf{\color{primary} Cột 4} & \textbf{\color{primary} Cột 5} & \textbf{\color{primary} Cột 6} & \textbf{\color{primary} Cột 7} & \textbf{\color{primary} Cột 8} & \textbf{\color{primary} Cột 9} & \textbf{\color{primary} Cột 10} \\
\midrule
\textbf{1.D} & \textbf{2.A} & \textbf{3.D} & \textbf{4.A} & \textbf{5.C} & \textbf{6.C} & \textbf{7.B} & \textbf{8.B} & \textbf{9.C} & \textbf{10.A} \\
\rowcolor{rowtint} \textbf{11.B} & \textbf{12.D} & \textbf{13.D} & \textbf{14.C} & \textbf{15.C} & \textbf{16.B} & \textbf{17.C} & \textbf{18.B} & \textbf{19.A} & \textbf{20.A} \\
\textbf{21.C} & \textbf{22.C} & \textbf{23.B} & \textbf{24.A} & \textbf{25.C} & \textbf{26.C} & \textbf{27.A} & \textbf{28.B} & \textbf{29.A} & \textbf{30.C} \\
\rowcolor{rowtint} \textbf{31.D} & \textbf{32.C} & \textbf{33.C} & \textbf{34.A} & \textbf{35.D} & \textbf{36.A} & \textbf{37.C} & \textbf{38.A} & \textbf{39.B} & \textbf{40.C} \\
\textbf{41.D} & \textbf{42.A} & \textbf{43.C} & \textbf{44.D} & \textbf{45.D} & \textbf{46.C} & \textbf{47.A} & \textbf{48.C} & \textbf{49.A} & \textbf{50.C} \\
\bottomrule
\end{tabularx}
\end{center}\vspace{10pt}
\section*{\color{primary} CHƯƠNG VI: XÁC SUẤT CÓ ĐIỀU KIỆN}
\textit{Trọng tâm kiến thức SGK: Xác suất có điều kiện P(A|B), công thức nhân xác suất, công thức xác suất toàn phần và công thức Bayes.}\par\vspace{4pt}
\subsection*{\color{accent} Phần I. Mức độ Dễ (Nhận biết -- Câu 1 đến Câu 20)}
\textbf{\color{primary} Câu 1.} Cho hai biến cố A và B với P(B) > 0. Công thức tính xác suất của biến cố A với điều kiện biến cố B đã xảy ra là:\par
\hspace*{1.5em}\textbf{A.} P(A|B) = P(A) $\cdot$ P(B)\quad\quad \textbf{B.} P(A|B) = P(AB) / P(B)\quad\quad \textbf{C.} P(A|B) = P(A) / P(B)\quad\quad \textbf{D.} P(A|B) = P(AB) / P(A)\par\vspace{2pt}
\textbf{\color{primary} Câu 2.} Thực hiện yêu cầu nhận biết sau: Nếu hai biến cố A và B độc lập với nhau và P(B) > 0 thì xác suất có điều kiện P(A|B) bằng.\par
\hspace*{1.5em}\textbf{A.} P(A|B) = P(B)\quad\quad \textbf{B.} P(A|B) = P(A)\quad\quad \textbf{C.} P(A|B) = 0\quad\quad \textbf{D.} P(A|B) = 1\par\vspace{2pt}
\textbf{\color{primary} Câu 3.} Trong bài học thuộc Chương VI, xét nội dung: Cho hai biến cố A và B thỏa mãn P(B) = 0,5 và P(AB) = 0,2. Khi đó xác suất có điều kiện P(A|B) bằng. Đáp án đúng là:\par
\hspace*{1.5em}\textbf{A.} 0,4\quad\quad \textbf{B.} 0,25\quad\quad \textbf{C.} 0,7\quad\quad \textbf{D.} 0,1\par\vspace{2pt}
\textbf{\color{primary} Câu 4.} Dựa vào kiến thức chuẩn SGK của Chương VI (Xác suất có điều kiện), cho hai biến cố A và B thỏa mãn P(B) = 6/14 và P(AB) = 2/14. Xác suất có điều kiện P(A|B) = P(AB)/P(B) bằng:\par
\hspace*{1.5em}\textbf{A.} 2/14\quad\quad \textbf{B.} 2/6\quad\quad \textbf{C.} 6/2\quad\quad \textbf{D.} 8/14\par\vspace{2pt}
\textbf{\color{primary} Câu 5.} Chọn kết quả đúng nhất khi xét: Cho hai biến cố A và B với P(B) > 0. Công thức tính xác suất của biến cố A với điều kiện biến cố B đã xảy ra là:\par
\hspace*{1.5em}\textbf{A.} P(A|B) = P(AB) / P(B)\quad\quad \textbf{B.} P(A|B) = P(A) $\cdot$ P(B)\quad\quad \textbf{C.} P(A|B) = P(AB) / P(A)\quad\quad \textbf{D.} P(A|B) = P(A) / P(B)\par\vspace{2pt}
\textbf{\color{primary} Câu 6.} Khi làm bài tập cơ bản về xác suất có điều kiện, một học sinh gặp câu hỏi: Nếu hai biến cố A và B độc lập với nhau và P(B) > 0 thì xác suất có điều kiện P(A|B) bằng. Kết quả nào dưới đây đúng?\par
\hspace*{1.5em}\textbf{A.} P(A|B) = P(B)\quad\quad \textbf{B.} P(A|B) = 0\quad\quad \textbf{C.} P(A|B) = 1\quad\quad \textbf{D.} P(A|B) = P(A)\par\vspace{2pt}
\textbf{\color{primary} Câu 7.} Xác định giá trị hoặc khẳng định đúng cho bài toán sau: Cho hai biến cố A và B thỏa mãn P(B) = 0,5 và P(AB) = 0,2. Khi đó xác suất có điều kiện P(A|B) bằng:\par
\hspace*{1.5em}\textbf{A.} 0,25\quad\quad \textbf{B.} 0,7\quad\quad \textbf{C.} 0,1\quad\quad \textbf{D.} 0,4\par\vspace{2pt}
\textbf{\color{primary} Câu 8.} Áp dụng trực tiếp định nghĩa và quy tắc trong Chương VI: Cho hai biến cố A và B thỏa mãn P(B) = 10/22 và P(AB) = 2/22. Xác suất có điều kiện P(A|B) = P(AB)/P(B) bằng:\par
\hspace*{1.5em}\textbf{A.} 10/2\quad\quad \textbf{B.} 12/22\quad\quad \textbf{C.} 2/22\quad\quad \textbf{D.} 2/10\par\vspace{2pt}
\textbf{\color{primary} Câu 9.} Kiểm tra nhanh kiến thức nền tảng: Cho hai biến cố A và B với P(B) > 0. Công thức tính xác suất của biến cố A với điều kiện biến cố B đã xảy ra là:\par
\hspace*{1.5em}\textbf{A.} P(A|B) = P(AB) / P(B)\quad\quad \textbf{B.} P(A|B) = P(A) $\cdot$ P(B)\quad\quad \textbf{C.} P(A|B) = P(A) / P(B)\quad\quad \textbf{D.} P(A|B) = P(AB) / P(A)\par\vspace{2pt}
\textbf{\color{primary} Câu 10.} Cho dữ kiện nhận biết sau đây: Nếu hai biến cố A và B độc lập với nhau và P(B) > 0 thì xác suất có điều kiện P(A|B) bằng. Phương án chính xác là:\par
\hspace*{1.5em}\textbf{A.} P(A|B) = P(B)\quad\quad \textbf{B.} P(A|B) = P(A)\quad\quad \textbf{C.} P(A|B) = 0\quad\quad \textbf{D.} P(A|B) = 1\par\vspace{2pt}
\textbf{\color{primary} Câu 11.} Hãy chỉ ra kết quả đúng trong bốn lựa chọn A, B, C, D đối với yêu cầu: Cho hai biến cố A và B thỏa mãn P(B) = 0,5 và P(AB) = 0,2. Khi đó xác suất có điều kiện P(A|B) bằng:\par
\hspace*{1.5em}\textbf{A.} 0,25\quad\quad \textbf{B.} 0,7\quad\quad \textbf{C.} 0,4\quad\quad \textbf{D.} 0,1\par\vspace{2pt}
\textbf{\color{primary} Câu 12.} Câu hỏi ôn tập lý thuyết và kỹ năng cơ bản (Chương VI): Cho hai biến cố A và B thỏa mãn P(B) = 14/30 và P(AB) = 2/30. Xác suất có điều kiện P(A|B) = P(AB)/P(B) bằng:\par
\hspace*{1.5em}\textbf{A.} 2/30\quad\quad \textbf{B.} 2/14\quad\quad \textbf{C.} 16/30\quad\quad \textbf{D.} 14/2\par\vspace{2pt}
\textbf{\color{primary} Câu 13.} Theo nội dung sách giáo khoa Kết nối tri thức với cuộc sống (Chương VI), cho hai biến cố A và B với P(B) > 0. Công thức tính xác suất của biến cố A với điều kiện biến cố B đã xảy ra là:\par
\hspace*{1.5em}\textbf{A.} P(A|B) = P(AB) / P(A)\quad\quad \textbf{B.} P(A|B) = P(A) / P(B)\quad\quad \textbf{C.} P(A|B) = P(A) $\cdot$ P(B)\quad\quad \textbf{D.} P(A|B) = P(AB) / P(B)\par\vspace{2pt}
\textbf{\color{primary} Câu 14.} Tìm đáp số đúng của bài toán cơ bản sau: Nếu hai biến cố A và B độc lập với nhau và P(B) > 0 thì xác suất có điều kiện P(A|B) bằng:\par
\hspace*{1.5em}\textbf{A.} P(A|B) = P(A)\quad\quad \textbf{B.} P(A|B) = 1\quad\quad \textbf{C.} P(A|B) = P(B)\quad\quad \textbf{D.} P(A|B) = 0\par\vspace{2pt}
\textbf{\color{primary} Câu 15.} Xét đối tượng toán học trong Chương VI: Cho hai biến cố A và B thỏa mãn P(B) = 0,5 và P(AB) = 0,2. Khi đó xác suất có điều kiện P(A|B) bằng. Giá trị hoặc tính chất tương ứng là:\par
\hspace*{1.5em}\textbf{A.} 0,1\quad\quad \textbf{B.} 0,25\quad\quad \textbf{C.} 0,7\quad\quad \textbf{D.} 0,4\par\vspace{2pt}
\textbf{\color{primary} Câu 16.} Hoàn thành phép xác định cơ bản sau đây: Cho hai biến cố A và B thỏa mãn P(B) = 18/38 và P(AB) = 2/38. Xác suất có điều kiện P(A|B) = P(AB)/P(B) bằng:\par
\hspace*{1.5em}\textbf{A.} 2/38\quad\quad \textbf{B.} 20/38\quad\quad \textbf{C.} 18/2\quad\quad \textbf{D.} 2/18\par\vspace{2pt}
\textbf{\color{primary} Câu 17.} Trong các khẳng định hoặc giá trị dưới đây, đâu là đáp án đúng cho câu hỏi: Cho hai biến cố A và B với P(B) > 0. Công thức tính xác suất của biến cố A với điều kiện biến cố B đã xảy ra là?\par
\hspace*{1.5em}\textbf{A.} P(A|B) = P(AB) / P(A)\quad\quad \textbf{B.} P(A|B) = P(A) $\cdot$ P(B)\quad\quad \textbf{C.} P(A|B) = P(AB) / P(B)\quad\quad \textbf{D.} P(A|B) = P(A) / P(B)\par\vspace{2pt}
\textbf{\color{primary} Câu 18.} Bài tập khởi động Chương VI: Nếu hai biến cố A và B độc lập với nhau và P(B) > 0 thì xác suất có điều kiện P(A|B) bằng:\par
\hspace*{1.5em}\textbf{A.} P(A|B) = 0\quad\quad \textbf{B.} P(A|B) = 1\quad\quad \textbf{C.} P(A|B) = P(B)\quad\quad \textbf{D.} P(A|B) = P(A)\par\vspace{2pt}
\textbf{\color{primary} Câu 19.} Quan sát dữ kiện đề bài: Cho hai biến cố A và B thỏa mãn P(B) = 0,5 và P(AB) = 0,2. Khi đó xác suất có điều kiện P(A|B) bằng. Hãy chọn phương án đúng:\par
\hspace*{1.5em}\textbf{A.} 0,7\quad\quad \textbf{B.} 0,4\quad\quad \textbf{C.} 0,1\quad\quad \textbf{D.} 0,25\par\vspace{2pt}
\textbf{\color{primary} Câu 20.} Sử dụng công thức hoặc tính chất đầu bài của Chương VI để giải quyết yêu cầu: Cho hai biến cố A và B thỏa mãn P(B) = 22/46 và P(AB) = 2/46. Xác suất có điều kiện P(A|B) = P(AB)/P(B) bằng:\par
\hspace*{1.5em}\textbf{A.} 24/46\quad\quad \textbf{B.} 2/22\quad\quad \textbf{C.} 22/2\quad\quad \textbf{D.} 2/46\par\vspace{2pt}
\subsection*{\color{accent} Phần II. Mức độ Trung bình (Thông hiểu -- Câu 21 đến Câu 40)}
\textbf{\color{primary} Câu 21.} Cho P(A) = 0,6 và P(B|A) = 0,5. Theo công thức nhân xác suất, giá trị của P(AB) bằng:\par
\hspace*{1.5em}\textbf{A.} 0,10\quad\quad \textbf{B.} 0,55\quad\quad \textbf{C.} 0,30\quad\quad \textbf{D.} 0,80\par\vspace{2pt}
\textbf{\color{primary} Câu 22.} Vận dụng quy tắc trong Chương VI (Xác suất có điều kiện) để giải bài toán: Cho P(B) = 0,4; P($\bar{B}$) = 0,6; P(A|B) = 0,5 và P(A|$\bar{B}$) = 0,2. Theo công thức xác suất toàn phần, giá trị của P(A) bằng:\par
\hspace*{1.5em}\textbf{A.} 0,20\quad\quad \textbf{B.} 0,70\quad\quad \textbf{C.} 0,32\quad\quad \textbf{D.} 0,12\par\vspace{2pt}
\textbf{\color{primary} Câu 23.} Biến đổi và tính toán theo kiến thức Chương VI: Một lớp học có 40\% học sinh nam và 60\% học sinh nữ. Biết rằng có 20\% học sinh nam cận thị và 30\% học sinh nữ cận thị. Chọn ngẫu nhiên 1 học sinh, xác suất để học sinh đó bị cận thị là. Kết quả thu được là:\par
\hspace*{1.5em}\textbf{A.} 0,18 (tức 18\%)\quad\quad \textbf{B.} 0,26 (tức 26\%)\quad\quad \textbf{C.} 0,50 (tức 50\%)\quad\quad \textbf{D.} 0,25 (tức 25\%)\par\vspace{2pt}
\textbf{\color{primary} Câu 24.} Xét bài toán mức độ thông hiểu sau: Cho hai biến cố A, B với P(A) = 3/9 và P(B|A) = 1/3. Theo công thức nhân xác suất, P(AB) = P(A)$\cdot$P(B|A) bằng. Đáp án đúng là:\par
\hspace*{1.5em}\textbf{A.} 1/9\quad\quad \textbf{B.} 4/9\quad\quad \textbf{C.} 1/27\quad\quad \textbf{D.} 3/9\par\vspace{2pt}
\textbf{\color{primary} Câu 25.} Để hoàn thành phiếu học tập Chương VI, hãy xác định: Cho P(A) = 0,6 và P(B|A) = 0,5. Theo công thức nhân xác suất, giá trị của P(AB) bằng:\par
\hspace*{1.5em}\textbf{A.} 0,10\quad\quad \textbf{B.} 0,30\quad\quad \textbf{C.} 0,80\quad\quad \textbf{D.} 0,55\par\vspace{2pt}
\textbf{\color{primary} Câu 26.} Phân tích dữ kiện bài toán: Cho P(B) = 0,4; P($\bar{B}$) = 0,6; P(A|B) = 0,5 và P(A|$\bar{B}$) = 0,2. Theo công thức xác suất toàn phần, giá trị của P(A) bằng. Giá trị cần tìm bằng:\par
\hspace*{1.5em}\textbf{A.} 0,12\quad\quad \textbf{B.} 0,32\quad\quad \textbf{C.} 0,70\quad\quad \textbf{D.} 0,20\par\vspace{2pt}
\textbf{\color{primary} Câu 27.} Áp dụng tính chất đã học ở Chương VI vào bài tập sau: Một lớp học có 40\% học sinh nam và 60\% học sinh nữ. Biết rằng có 20\% học sinh nam cận thị và 30\% học sinh nữ cận thị. Chọn ngẫu nhiên 1 học sinh, xác suất để học sinh đó bị cận thị là:\par
\hspace*{1.5em}\textbf{A.} 0,25 (tức 25\%)\quad\quad \textbf{B.} 0,26 (tức 26\%)\quad\quad \textbf{C.} 0,18 (tức 18\%)\quad\quad \textbf{D.} 0,50 (tức 50\%)\par\vspace{2pt}
\textbf{\color{primary} Câu 28.} Thực hiện các bước suy luận hợp lý để giải quyết yêu cầu: Cho hai biến cố A, B với P(A) = 3/13 và P(B|A) = 1/3. Theo công thức nhân xác suất, P(AB) = P(A)$\cdot$P(B|A) bằng:\par
\hspace*{1.5em}\textbf{A.} 4/13\quad\quad \textbf{B.} 1/39\quad\quad \textbf{C.} 1/13\quad\quad \textbf{D.} 3/13\par\vspace{2pt}
\textbf{\color{primary} Câu 29.} Cho bài toán tự luyện (Xác suất có điều kiện): Cho P(A) = 0,6 và P(B|A) = 0,5. Theo công thức nhân xác suất, giá trị của P(AB) bằng. Phương án nào sau đây chính xác?\par
\hspace*{1.5em}\textbf{A.} 0,10\quad\quad \textbf{B.} 0,30\quad\quad \textbf{C.} 0,80\quad\quad \textbf{D.} 0,55\par\vspace{2pt}
\textbf{\color{primary} Câu 30.} Kết hợp các phép biến đổi trong Chương VI, hãy tìm kết quả của bài toán: Cho P(B) = 0,4; P($\bar{B}$) = 0,6; P(A|B) = 0,5 và P(A|$\bar{B}$) = 0,2. Theo công thức xác suất toàn phần, giá trị của P(A) bằng:\par
\hspace*{1.5em}\textbf{A.} 0,12\quad\quad \textbf{B.} 0,20\quad\quad \textbf{C.} 0,32\quad\quad \textbf{D.} 0,70\par\vspace{2pt}
\textbf{\color{primary} Câu 31.} Một bài kiểm tra định kỳ nội dung Chương VI nêu câu hỏi: Một lớp học có 40\% học sinh nam và 60\% học sinh nữ. Biết rằng có 20\% học sinh nam cận thị và 30\% học sinh nữ cận thị. Chọn ngẫu nhiên 1 học sinh, xác suất để học sinh đó bị cận thị là. Hãy chọn đáp số đúng:\par
\hspace*{1.5em}\textbf{A.} 0,18 (tức 18\%)\quad\quad \textbf{B.} 0,26 (tức 26\%)\quad\quad \textbf{C.} 0,50 (tức 50\%)\quad\quad \textbf{D.} 0,25 (tức 25\%)\par\vspace{2pt}
\textbf{\color{primary} Câu 32.} Tìm lời giải đúng cho bài toán thông hiểu: Cho hai biến cố A, B với P(A) = 3/17 và P(B|A) = 1/3. Theo công thức nhân xác suất, P(AB) = P(A)$\cdot$P(B|A) bằng:\par
\hspace*{1.5em}\textbf{A.} 1/51\quad\quad \textbf{B.} 1/17\quad\quad \textbf{C.} 4/17\quad\quad \textbf{D.} 3/17\par\vspace{2pt}
\textbf{\color{primary} Câu 33.} Sử dụng mối liên hệ giữa các đại lượng trong Chương VI: Cho P(A) = 0,6 và P(B|A) = 0,5. Theo công thức nhân xác suất, giá trị của P(AB) bằng:\par
\hspace*{1.5em}\textbf{A.} 0,30\quad\quad \textbf{B.} 0,10\quad\quad \textbf{C.} 0,55\quad\quad \textbf{D.} 0,80\par\vspace{2pt}
\textbf{\color{primary} Câu 34.} Hãy tính toán và chọn phương án đúng đối với bài tập: Cho P(B) = 0,4; P($\bar{B}$) = 0,6; P(A|B) = 0,5 và P(A|$\bar{B}$) = 0,2. Theo công thức xác suất toàn phần, giá trị của P(A) bằng:\par
\hspace*{1.5em}\textbf{A.} 0,12\quad\quad \textbf{B.} 0,20\quad\quad \textbf{C.} 0,32\quad\quad \textbf{D.} 0,70\par\vspace{2pt}
\textbf{\color{primary} Câu 35.} Dựa trên các tính chất trọng tâm của Chương VI, xét bài toán: Một lớp học có 40\% học sinh nam và 60\% học sinh nữ. Biết rằng có 20\% học sinh nam cận thị và 30\% học sinh nữ cận thị. Chọn ngẫu nhiên 1 học sinh, xác suất để học sinh đó bị cận thị là:\par
\hspace*{1.5em}\textbf{A.} 0,26 (tức 26\%)\quad\quad \textbf{B.} 0,25 (tức 25\%)\quad\quad \textbf{C.} 0,50 (tức 50\%)\quad\quad \textbf{D.} 0,18 (tức 18\%)\par\vspace{2pt}
\textbf{\color{primary} Câu 36.} Thực hành giải toán (Xác suất có điều kiện): Cho hai biến cố A, B với P(A) = 3/21 và P(B|A) = 1/3. Theo công thức nhân xác suất, P(AB) = P(A)$\cdot$P(B|A) bằng. Đáp án nào dưới đây thỏa mãn?\par
\hspace*{1.5em}\textbf{A.} 3/21\quad\quad \textbf{B.} 1/21\quad\quad \textbf{C.} 1/63\quad\quad \textbf{D.} 4/21\par\vspace{2pt}
\textbf{\color{primary} Câu 37.} Kiểm tra kỹ năng biến đổi và suy luận qua bài toán: Cho P(A) = 0,6 và P(B|A) = 0,5. Theo công thức nhân xác suất, giá trị của P(AB) bằng:\par
\hspace*{1.5em}\textbf{A.} 0,30\quad\quad \textbf{B.} 0,80\quad\quad \textbf{C.} 0,10\quad\quad \textbf{D.} 0,55\par\vspace{2pt}
\textbf{\color{primary} Câu 38.} Cho giả thiết thuộc phạm vi Chương VI: Cho P(B) = 0,4; P($\bar{B}$) = 0,6; P(A|B) = 0,5 và P(A|$\bar{B}$) = 0,2. Theo công thức xác suất toàn phần, giá trị của P(A) bằng. Kết luận đúng là:\par
\hspace*{1.5em}\textbf{A.} 0,32\quad\quad \textbf{B.} 0,12\quad\quad \textbf{C.} 0,20\quad\quad \textbf{D.} 0,70\par\vspace{2pt}
\textbf{\color{primary} Câu 39.} Rút gọn hoặc tính giá trị theo yêu cầu sau của Chương VI: Một lớp học có 40\% học sinh nam và 60\% học sinh nữ. Biết rằng có 20\% học sinh nam cận thị và 30\% học sinh nữ cận thị. Chọn ngẫu nhiên 1 học sinh, xác suất để học sinh đó bị cận thị là:\par
\hspace*{1.5em}\textbf{A.} 0,50 (tức 50\%)\quad\quad \textbf{B.} 0,25 (tức 25\%)\quad\quad \textbf{C.} 0,18 (tức 18\%)\quad\quad \textbf{D.} 0,26 (tức 26\%)\par\vspace{2pt}
\textbf{\color{primary} Câu 40.} Xác định phương án đúng sau khi thực hiện phép tính / suy luận: Cho hai biến cố A, B với P(A) = 3/25 và P(B|A) = 1/3. Theo công thức nhân xác suất, P(AB) = P(A)$\cdot$P(B|A) bằng:\par
\hspace*{1.5em}\textbf{A.} 1/75\quad\quad \textbf{B.} 1/25\quad\quad \textbf{C.} 4/25\quad\quad \textbf{D.} 3/25\par\vspace{2pt}
\subsection*{\color{accent} Phần III. Mức độ Khó (Vận dụng \& Vận dụng cao -- Câu 41 đến Câu 50)}
\textbf{\color{primary} Câu 41.} Một nhà máy có hai phân xưởng I và II lần lượt sản xuất 60\% và 40\% tổng sản phẩm. Tỉ lệ phế phẩm của phân xưởng I là 2\% và của phân xưởng II là 3\%. Lấy ngẫu nhiên 1 sản phẩm thấy đó là phế phẩm, theo công thức Bayes xác suất để phế phẩm đó do phân xưởng I sản xuất bằng:\par
\hspace*{1.5em}\textbf{A.} 2/3\quad\quad \textbf{B.} 2/5 (tức 40\%)\quad\quad \textbf{C.} 3/5 (tức 60\%)\quad\quad \textbf{D.} 1/2 (tức 50\%)\par\vspace{2pt}
\textbf{\color{primary} Câu 42.} Bài toán vận dụng (Chương VI – Xác suất có điều kiện): Tỉ lệ người mắc một bệnh X trong cộng đồng là 1\%. Một xét nghiệm cho kết quả dương tính với 95\% người mắc bệnh và cho dương tính giả với 2\% người không mắc bệnh. Theo công thức Bayes, nếu một người có kết quả xét nghiệm dương tính thì xác suất người đó thực sự mắc bệnh X là. Kết quả đúng là:\par
\hspace*{1.5em}\textbf{A.} 2 / 293\quad\quad \textbf{B.} 95 / 293 ($\approx$ 32,4\%)\quad\quad \textbf{C.} 19 / 20\quad\quad \textbf{D.} 95\%\par\vspace{2pt}
\textbf{\color{primary} Câu 43.} Vận dụng kiến thức nâng cao trong phạm vi Chương VI để giải bài toán: Một hộp có 2 loại sản phẩm: phân xưởng I chiếm 4/9 tổng số sản phẩm (tỉ lệ phế phẩm 2\%) và phân xưởng II chiếm 5/9 tổng số sản phẩm (tỉ lệ phế phẩm 2\%). Chọn ngẫu nhiên 1 sản phẩm, theo công thức xác suất toàn phần thì xác suất chọn phải phế phẩm là:\par
\hspace*{1.5em}\textbf{A.} 8/9\%\quad\quad \textbf{B.} 0,02 (tức 2\%)\quad\quad \textbf{C.} 0,01 (tức 1\%)\quad\quad \textbf{D.} 0,04 (tức 4\%)\par\vspace{2pt}
\textbf{\color{primary} Câu 44.} Phân tích và giải bài toán tổng hợp sau của Chương VI: Một nhà máy có hai phân xưởng I và II lần lượt sản xuất 60\% và 40\% tổng sản phẩm. Tỉ lệ phế phẩm của phân xưởng I là 2\% và của phân xưởng II là 3\%. Lấy ngẫu nhiên 1 sản phẩm thấy đó là phế phẩm, theo công thức Bayes xác suất để phế phẩm đó do phân xưởng I sản xuất bằng:\par
\hspace*{1.5em}\textbf{A.} 1/2 (tức 50\%)\quad\quad \textbf{B.} 3/5 (tức 60\%)\quad\quad \textbf{C.} 2/5 (tức 40\%)\quad\quad \textbf{D.} 2/3\par\vspace{2pt}
\textbf{\color{primary} Câu 45.} Thử thách tư duy trong phạm vi kiến thức Chương VI: Tỉ lệ người mắc một bệnh X trong cộng đồng là 1\%. Một xét nghiệm cho kết quả dương tính với 95\% người mắc bệnh và cho dương tính giả với 2\% người không mắc bệnh. Theo công thức Bayes, nếu một người có kết quả xét nghiệm dương tính thì xác suất người đó thực sự mắc bệnh X là. Đáp số chính xác là:\par
\hspace*{1.5em}\textbf{A.} 19 / 20\quad\quad \textbf{B.} 95 / 293 ($\approx$ 32,4\%)\quad\quad \textbf{C.} 2 / 293\quad\quad \textbf{D.} 95\%\par\vspace{2pt}
\textbf{\color{primary} Câu 46.} Sử dụng phương pháp giải toán của Chương VI để tìm đáp án cho bài toán: Một hộp có 2 loại sản phẩm: phân xưởng I chiếm 7/15 tổng số sản phẩm (tỉ lệ phế phẩm 2\%) và phân xưởng II chiếm 8/15 tổng số sản phẩm (tỉ lệ phế phẩm 2\%). Chọn ngẫu nhiên 1 sản phẩm, theo công thức xác suất toàn phần thì xác suất chọn phải phế phẩm là:\par
\hspace*{1.5em}\textbf{A.} 0,02 (tức 2\%)\quad\quad \textbf{B.} 14/15\%\quad\quad \textbf{C.} 0,01 (tức 1\%)\quad\quad \textbf{D.} 0,04 (tức 4\%)\par\vspace{2pt}
\textbf{\color{primary} Câu 47.} Xét bài toán vận dụng thực tiễn / nâng cao thuộc Chương VI: Một nhà máy có hai phân xưởng I và II lần lượt sản xuất 60\% và 40\% tổng sản phẩm. Tỉ lệ phế phẩm của phân xưởng I là 2\% và của phân xưởng II là 3\%. Lấy ngẫu nhiên 1 sản phẩm thấy đó là phế phẩm, theo công thức Bayes xác suất để phế phẩm đó do phân xưởng I sản xuất bằng:\par
\hspace*{1.5em}\textbf{A.} 1/2 (tức 50\%)\quad\quad \textbf{B.} 2/3\quad\quad \textbf{C.} 2/5 (tức 40\%)\quad\quad \textbf{D.} 3/5 (tức 60\%)\par\vspace{2pt}
\textbf{\color{primary} Câu 48.} Tìm kết quả cuối cùng của bài toán phân loại học sinh khá giỏi (Chương VI): Tỉ lệ người mắc một bệnh X trong cộng đồng là 1\%. Một xét nghiệm cho kết quả dương tính với 95\% người mắc bệnh và cho dương tính giả với 2\% người không mắc bệnh. Theo công thức Bayes, nếu một người có kết quả xét nghiệm dương tính thì xác suất người đó thực sự mắc bệnh X là:\par
\hspace*{1.5em}\textbf{A.} 95 / 293 ($\approx$ 32,4\%)\quad\quad \textbf{B.} 2 / 293\quad\quad \textbf{C.} 95\%\quad\quad \textbf{D.} 19 / 20\par\vspace{2pt}
\textbf{\color{primary} Câu 49.} Kết hợp các tính chất đã học trong Chương VI, hãy giải quyết bài toán: Một hộp có 2 loại sản phẩm: phân xưởng I chiếm 10/21 tổng số sản phẩm (tỉ lệ phế phẩm 2\%) và phân xưởng II chiếm 11/21 tổng số sản phẩm (tỉ lệ phế phẩm 2\%). Chọn ngẫu nhiên 1 sản phẩm, theo công thức xác suất toàn phần thì xác suất chọn phải phế phẩm là:\par
\hspace*{1.5em}\textbf{A.} 0,01 (tức 1\%)\quad\quad \textbf{B.} 0,02 (tức 2\%)\quad\quad \textbf{C.} 20/21\%\quad\quad \textbf{D.} 0,04 (tức 4\%)\par\vspace{2pt}
\textbf{\color{primary} Câu 50.} Câu hỏi vận dụng cao cuối Chương VI: Một nhà máy có hai phân xưởng I và II lần lượt sản xuất 60\% và 40\% tổng sản phẩm. Tỉ lệ phế phẩm của phân xưởng I là 2\% và của phân xưởng II là 3\%. Lấy ngẫu nhiên 1 sản phẩm thấy đó là phế phẩm, theo công thức Bayes xác suất để phế phẩm đó do phân xưởng I sản xuất bằng. Phương án đúng là:\par
\hspace*{1.5em}\textbf{A.} 2/3\quad\quad \textbf{B.} 1/2 (tức 50\%)\quad\quad \textbf{C.} 2/5 (tức 40\%)\quad\quad \textbf{D.} 3/5 (tức 60\%)\par\vspace{2pt}
\subsection*{\color{primary} Bảng đáp án Chương VI: Xác suất có điều kiện}
\begin{center}
\renewcommand{\arraystretch}{1.25}
\begin{tabularx}{\textwidth}{*{10}{>{\centering\arraybackslash}X}}
\toprule
\textbf{\color{primary} Cột 1} & \textbf{\color{primary} Cột 2} & \textbf{\color{primary} Cột 3} & \textbf{\color{primary} Cột 4} & \textbf{\color{primary} Cột 5} & \textbf{\color{primary} Cột 6} & \textbf{\color{primary} Cột 7} & \textbf{\color{primary} Cột 8} & \textbf{\color{primary} Cột 9} & \textbf{\color{primary} Cột 10} \\
\midrule
\textbf{1.B} & \textbf{2.B} & \textbf{3.A} & \textbf{4.B} & \textbf{5.A} & \textbf{6.D} & \textbf{7.D} & \textbf{8.D} & \textbf{9.A} & \textbf{10.B} \\
\rowcolor{rowtint} \textbf{11.C} & \textbf{12.B} & \textbf{13.D} & \textbf{14.A} & \textbf{15.D} & \textbf{16.D} & \textbf{17.C} & \textbf{18.D} & \textbf{19.B} & \textbf{20.B} \\
\textbf{21.C} & \textbf{22.C} & \textbf{23.B} & \textbf{24.A} & \textbf{25.B} & \textbf{26.B} & \textbf{27.B} & \textbf{28.C} & \textbf{29.B} & \textbf{30.C} \\
\rowcolor{rowtint} \textbf{31.B} & \textbf{32.B} & \textbf{33.A} & \textbf{34.C} & \textbf{35.A} & \textbf{36.B} & \textbf{37.A} & \textbf{38.A} & \textbf{39.D} & \textbf{40.B} \\
\textbf{41.D} & \textbf{42.B} & \textbf{43.B} & \textbf{44.A} & \textbf{45.B} & \textbf{46.A} & \textbf{47.A} & \textbf{48.A} & \textbf{49.B} & \textbf{50.B} \\
\bottomrule
\end{tabularx}
\end{center}\vspace{10pt}
\end{document}
